\subsection{The Hall--Connor nonexistence theorem}\label{hc:sec:main}
\begin{mergedproofsection}
\raggedbottom

We will follow the ordinary proof of
\path{cgss/proofs/hallconnor/hallconnor.tex}. The argument proves the
nonexistence assertion used in Remark 5.12 and Example 5.13 of
\cite{CGSS}, together with its complementary parameters. We use
Hall and Connor's argument \cite{HC}, including proofs of the instances
they take from Connor. We have not consulted Connor's paper directly.
We omit the two comparison remarks of the ordinary proof, which the
argument does not use.

We follow the blueprint conventions and Lamport \cite{L95}:
\textsc{Assume} introduces hypotheses, \textsc{Prove} the conclusion,
\textsc{Let} a definition, and \textsc{Pick} a witness.
Each numbered step has a proof. A reference $\langle0\rangle{:}i$
means hypothesis $i$ of the result; $\langle k\rangle$ means the
assumption of the enclosing level-$k$ step. Closing \textsc{Qed}
steps state why the preceding steps suffice. The tags \new\ and
\repair\ mark a change of argument and a filled gap, respectively.
The no-op commands \verb|\lean| and \verb|\uses| record the Lean
declarations and the direct dependencies.

We use the notation of \path{cgss/notation.md}, Sections 7 and 10:
\texttt{Design P B v k b r lam}, \path{D.incMatrix}, and \path{D.compl}.
The Hall--Connor names not listed individually there retain the names
of the ordinary proof, following the namespace and
naming rules of Section 10. Points and block labels are separate finite
types with decidable equality. All sets used in counts are finite sets;
their cardinalities are natural numbers. Moment identities and
subtractions in weighted sums are evaluated in $\ints$ or $\rats$,
never by truncated natural subtraction. Matrix entries are $M(u,w)$.
We write $I_U$ for the identity, $J_{U,V}$ for the all-ones matrix,
and $\Zj_U$ for the all-ones vector. Omitted index types are determined
by the products. All ranks are over $\rats$.

\subsubsection{Definitions and Mathlib inputs}

\begin{namedresult}{Definition}{(labelled designs)}\label{hc:def:design}
 \lean{Design}\lean{Design.IsSymmetric}\uses{}
 \leavevmode\par
 \assume{\begin{pfenum}
  \item $P,\cB$ are finite types with decidable equality.
  \item $v,k,b,r,\lambda\in\nn$, with $v>1$, $k>1$, and $b>0$.
 \end{pfenum}}
 A design $D$ with these parameters consists of a function
 $B\mapsto D(B)\subseteq P$ such that
 \[
 \begin{gathered}
 |P|=v,\quad |\cB|=b,\quad |D(B)|=k\quad(B\in\cB),\\
 |\{B:p\in D(B)\}|=r\quad(p\in P),\\
 |\{B:p,q\in D(B)\}|=\lambda\quad(p,q\in P,\ p\ne q).
 \end{gathered}
 \]
 We write $p\in B$ for $p\in D(B)$; the map $D$ need not be injective.
 The predicate $D.\mathrm{IsSymmetric}$ means $v=b$.
 In Lean the parameter names are \texttt{v k b r lam}; any positivity
 fields of the shared interface are satisfied by the parameters here.
\end{namedresult}

\begin{namedresult}{Definition}{(incidence and block matrices)}\label{hc:def:matrices}
 \lean{Design.incMatrix}\lean{Design.G}\lean{Design.C}\uses{hc:def:design}
 \leavevmode\par
 \assume{\begin{pfenum}
  \item $P,\cB$ are finite types with decidable equality.
  \item $D$ is a design with parameters $(v,k,b,r,\lambda)$.
 \end{pfenum}}
 \pflet{$N\triangleq D.\mathrm{incMatrix}\in\mat{P}{\cB}{\ints}$, where}
 \[
 N(p,B)\triangleq\begin{cases}1,&p\in B,\\0,&p\notin B.\end{cases}
 \]
 Write $N_{\rats}$ for its entrywise rational cast. Put
 $G\triangleq N_{\rats}^TN_{\rats}$ and
 $C\triangleq35I_{\cB}-7G+10J_{\cB}$, both over $\rats$.
 A product containing $G$ or $C$ uses $N_{\rats}$; we abbreviate it
 by $N$ in rational calculations. Products asserted integral use the
 original $N$.
\end{namedresult}

\begin{namedresult}{Definition}{(the block graph and all six-cliques)}\label{hc:def:graph}
 \lean{Design.Z}\lean{Design.sixCliques}\uses{hc:def:design}
 \leavevmode\par
 \assume{\begin{pfenum}
  \item $P,\cB$ are finite types with decidable equality.
  \item $D$ is a design with parameters $(v,k,b,r,\lambda)$.
 \end{pfenum}}
 \pflet{$Z$ be the simple graph on $\cB$ with}
 \[
 B\sim B'\quad\Longleftrightarrow\quad B\ne B'\ \text{and}\ |B\cap B'|=1.
 \]
 Here $B\cap B'$ means $D(B)\cap D(B')$.
 Put $\Gamma(B)\triangleq\{B':B\sim B'\}$ and
 \[
 \cS\triangleq D.\mathrm{sixCliques}
   \triangleq\{S\subseteq\cB:|S|=6\ \text{and}\
      \forall X,Y\in S,\ X\ne Y\Rightarrow X\sim Y\}.
 \]
 Thus $\cS$ is the finite set of \emph{all} six-cliques, not a choice
 of local cliques. Its members form a finite subtype when used as indices.
\end{namedresult}

\begingroup\raggedright
\paragraph{Finite counting and order.}\factlabel{F1}{hc:facts:finite}
\uses{}
All declarations below were checked in Mathlib at commit
\path{082e2d37} (4 September 2026), the version pinned in
\path{lake-manifest.json}; paths are relative to \path{Mathlib/}. Finite double counting is
interchange of indicator sums, followed by grouping fibres:
\begin{itemize}
 \item \path{Finset.sum_comm}, in
 \path{Algebra/BigOperators/Group/Finset/Sigma.lean};
 \item \path{Finset.card_eq_sum_card_fiberwise}, in
 \path{Algebra/BigOperators/Group/Finset/Basic.lean};
 \item \path{Finset.card_sdiff_of_subset},
 \path{Finset.card_union_add_card_inter},
 \path{Finset.eq_of_subset_of_card_le}, \path{Finset.card_pos}, and
 \path{Finset.card_le_one}, \path{Finset.card_eq_one},
 \path{Finset.card_eq_two}, \path{Finset.card_eq_three}, and
 \path{Finset.card_image_of_injective}, in \path{Data/Finset/Card.lean};
 \item \path{Fintype.bijective_iff_injective_and_card} and
 \path{Fintype.equivFin}, in \path{Data/Fintype/EquivFin.lean};
 \item \path{Fintype.card_sum}, in \path{Data/Fintype/Sum.lean},
 and \path{Fintype.card_option}, in \path{Data/Fintype/Option.lean};
 \item \path{Fintype.sum_sum_type} and \path{Fintype.sum_option}, in
 \path{Data/Fintype/BigOperators.lean};
 \item \path{Finset.sum_eq_zero_iff_of_nonneg}, in
 \path{Algebra/Order/BigOperators/Group/Finset.lean}.
\end{itemize}
The last declaration says that a sum of nonnegative terms is zero
exactly when every term is zero. The equal-cardinality theorem turns
an injection between finite types of equal size into a bijection.
An enumeration of a finite subset is an equivalence to its subtype;
its inclusion in the ambient type is injective. Integer and rational
identities below use ring arithmetic, nonnegative squares, and finite
sum distributivity. Small numerical tables are explicit computations.

\paragraph{Matrix operations and reindexing.}\factlabel{F2}{hc:facts:matrix}
\uses{}
We use \path{Matrix.mul_apply}, \path{Matrix.mul_assoc},
\path{Matrix.transpose_mul}, \path{Matrix.dotProduct_mulVec},
\path{Matrix.mulVec_mulVec}, \path{Matrix.map_mul}, and
\path{Matrix.submatrix_mul_equiv}, all in \path{Data/Matrix/Mul.lean}.
The last identity says that, for
an equivalence $e$ of the summed index type,
\[
 (A\mathbin{\mathrm{submatrix}}f\,e)
 (B\mathbin{\mathrm{submatrix}}e\,g)
       =(AB)\mathbin{\mathrm{submatrix}}f\,g.
\]
We also use \path{Matrix.transpose_map} and
\path{Matrix.transpose_submatrix}, in
\path{LinearAlgebra/Matrix/Defs.lean}, and \path{Matrix.map_one}, in
\path{Data/Matrix/Diagonal.lean}. Applied to the ring homomorphism
$\ints\to\rats$, these transport products, transposes, and $I$.
The cast of $J$ is $J$ entrywise because $1$ maps to $1$; addition
and scalar multiplication commute with the cast entrywise.

\paragraph{Ranks and determinants.}\factlabel{F3}{hc:facts:rank}
\uses{}
The following declarations are in
\path{LinearAlgebra/Matrix/Rank.lean}:
\path{Matrix.rank_self_mul_transpose}, \path{Matrix.rank_transpose},
\path{Matrix.rank_of_isUnit}, \path{Matrix.rank_of_det_ne_zero},
\path{Matrix.rank_add_rank_le_card_of_mul_eq_zero}, and
\path{Matrix.rank_submatrix_le}.
The identity $\operatorname{rk}(NN^T)=\operatorname{rk}N$ applies over
the linearly ordered field $\rats$; rank under transpose applies over
a field. A nonzero determinant gives full rank. If $AB=0$, the sum
of their ranks is at most the size of the shared index type.
\path{Matrix.isUnit_iff_isUnit_det}, in
\path{LinearAlgebra/Matrix/NonsingularInverse.lean}, turns a nonzero
rational determinant into an invertible matrix.

\paragraph{Determinant arithmetic.}\factlabel{F4}{hc:facts:det}
\uses{}
We use \path{Matrix.det_mul}, \path{Matrix.det_transpose},
\path{Matrix.det_smul}, \path{Matrix.det_one}, and \path{Int.cast_det}, in
\path{LinearAlgebra/Matrix/Determinant/Basic.lean}.
We use the matrix determinant lemma
\path{Matrix.det_one_add_replicateCol_mul_replicateRow}, in
\path{LinearAlgebra/Matrix/SchurComplement.lean}, in the form
$\det(I+\Zu\Zv^T)=1+\Zv^T\Zu$ (the shared singleton index is
\texttt{Unit}). These give multiplicativity, scalar homogeneity,
and compatibility with the cast $\ints\to\rats$. In particular, the
determinant over $\rats$ of an integral matrix is the cast of its
integer determinant. We do not use the analytic
positive-semidefinite determinant theorem, whose scalar hypotheses
do not include $\rats$.

\paragraph{Integer divisibility.}\factlabel{F5}{hc:facts:prime}
\uses{}
We use \path{Prime.dvd_of_dvd_pow}, in \path{Algebra/Prime/Defs.lean},
\path{Int.prime_iff_natAbs_prime}, in \path{RingTheory/Int/Basic.lean},
and \path{Nat.prime_five}, in \path{Data/Nat/Prime/Defs.lean}.
The last two prove that $5$ is prime in $\ints$; the first gives
$5\mid d^2\Rightarrow5\mid d$ for $d\in\ints$.

\par\endgroup

\subsubsection{Counts and intersections}

\begin{namedresult}{Lemma}{(the forced parameters)}\label{hc:lem:parameters}
 \lean{Design.exists_of_pair_counts_15_5_2}
 \uses{hc:def:design,hc:facts:finite}
 \leavevmode\par
 \assume{\begin{pfenum}
  \item $P,\cB$ are finite types with decidable equality and $|P|=15$.
  \item $F:\cB\to\operatorname{Finset}(P)$ has $|F(B)|=5$ for every $B$.
  \item Each pair of distinct points lies in exactly two labelled blocks of $F$.
 \end{pfenum}}
 \prove{Every point is in seven labels, $|\cB|=21$, and
 \[
 \exists D:\texttt{Design P B 15 5 21 7 2},\
       \forall B\in\cB,\ D.\mathrm{block}(B)=F(B).
 \]
 Here the Lean type variable \texttt{B} denotes $\cB$.}
\end{namedresult}
\begin{proof}
 \pfsketch\ We will count first through one point, then over all points.
 \step{hcParPoint}{\assume{$p\in P$.}
  \prove{$r_p\triangleq|\{B:p\in F(B)\}|=7$.}}
 \begin{proof}
  \step{hcParBlocks}{The set of pairs $(q,B)$ with $q\ne p$ and
   $p,q\in F(B)$ has size $4r_p$.}
  \begin{proof}\pf\ For each of the $r_p$ labels, deleting $p$ from
   its five points leaves four choices, by \asm[2]{0} and
   the fibre count in \ref{hc:facts:finite}.~\qed\end{proof}
  \step{hcParPairs}{The same set has size $28$.}
  \begin{proof}\pf\ There are $15-1=14$ choices of $q$ and two labels
   for each, by \asm[1]{0}, \asm[3]{0} and
   \ref{hc:facts:finite}.~\qed\end{proof}
  \qedstep
  \begin{proof}\pf\ The equality $4r_p=28$ from
   \stepref{hcParBlocks} and \stepref{hcParPairs} gives $r_p=7$.~\qed\end{proof}
 \end{proof}
 \step{hcParTotal}{$5|\cB|=105$ and $|\cB|=21$.}
 \begin{proof}\pf\ Count $(p,B)$ with $p\in F(B)$ by blocks and by
  points, using \asm[1]{0}, \asm[2]{0}, \stepref{hcParPoint}, and
  \ref{hc:facts:finite}: $5|\cB|=15\cdot7$. Division by five gives
  the second assertion.~\qed\end{proof}
 \qedstep
 \begin{proof}\pf\ Steps \stepref{hcParPoint} and \stepref{hcParTotal},
  with \asm[1]{0}, \asm[2]{0}, and \asm[3]{0}, give every field of
  Definition~\ref{hc:def:design}. Choose the structure with block
  function $F$; its block equality is reflexivity, and
  $15>1$, $5>1$, $21>0$ are numerical.~\qed\end{proof}
\end{proof}

\begin{namedresult}{Lemma}{(incidence identities)}\label{hc:lem:incidence}
 \lean{Design.incMatrix_mul_transpose_15_5_2}
 \lean{Design.incMatrix_row_sum_15_5_2}\lean{Design.incMatrix_col_sum_15_5_2}
 \lean{Design.G_apply}\lean{Design.G_mul_G}\lean{Design.G_mul_J}
 \uses{hc:def:design,hc:def:matrices,hc:facts:finite,hc:facts:matrix}
 \leavevmode\par
 \assume{\begin{pfenum}
  \item $P,\cB$ are finite types with decidable equality.
  \item $D$ is a design with parameters $(15,5,21,7,2)$.
 \end{pfenum}}
 \prove{The integral incidence matrix has
 \[
 NN^T=5I_P+2J_P,\quad N\Zj_{\cB}=7\Zj_P,\quad
 N^T\Zj_P=5\Zj_{\cB}.
 \]
 Over $\rats$, $G(B,B')=|B\cap B'|$, $G^T=G$, and
 \[
 G^2=5G+50J_{\cB},\quad GJ_{\cB}=J_{\cB}G=35J_{\cB},\quad
 J_{\cB}^2=21J_{\cB}.
 \]}
\end{namedresult}
\begin{proof}
 \pfsketch\ The matrix entries count points or labels; associativity
 then gives the identities for $G$.
 \step{hcIncDiag}{$(NN^T)(p,p)=7$ for every $p\in P$.}
 \begin{proof}\pf\ By Definitions~\ref{hc:def:matrices} and
  \ref{hc:def:design}, the entry is the sum of the squares of the
  seven incidence indicators equal to one; use \ref{hc:facts:matrix}.~\qed\end{proof}
 \step{hcIncOff}{\assume{$p,q\in P$ and $p\ne q$.}
  \prove{$(NN^T)(p,q)=2$.}}
 \begin{proof}\pf\ The product of the two indicators is one exactly
  on labels containing $p,q$, of which there are two by
  Definition~\ref{hc:def:design}; expand by \ref{hc:facts:matrix}.~\qed\end{proof}
 \step{hcIncNN}{$NN^T=5I_P+2J_P$.}
 \begin{proof}\pf\ Steps \stepref{hcIncDiag} and \stepref{hcIncOff}
  treat the exhaustive alternatives $p=q$ and $p\ne q$.~\qed\end{proof}
 \step{hcIncRow}{$N\Zj_{\cB}=7\Zj_P$.}
 \begin{proof}\pf\ The row sum counts the seven labels on $p$,
  by Definition~\ref{hc:def:design} and \ref{hc:facts:matrix}.~\qed\end{proof}
 \step{hcIncCol}{$N^T\Zj_P=5\Zj_{\cB}$.}
 \begin{proof}\pf\ The column sum counts the five points of a block,
  by Definition~\ref{hc:def:design} and \ref{hc:facts:matrix}.~\qed\end{proof}
 \step{hcIncG}{$G(B,B')=|B\cap B'|$ and $G^T=G$.}
 \begin{proof}\pf\ In Definition~\ref{hc:def:matrices}, the summand
  $N(p,B)N(p,B')$ is the intersection indicator. Its symmetry in
  $B,B'$ gives symmetry of $G$, by \ref{hc:facts:matrix}.~\qed\end{proof}
 \step{hcIncSquare}{$G^2=5G+50J_{\cB}$.}
 \begin{proof}
  \pf\ Cast \stepref{hcIncNN} and \stepref{hcIncCol} entrywise
  using the product, transpose, and constant-matrix rules of
  \ref{hc:facts:matrix}. Then
  \begin{align*}
   G^2&=N^T(NN^T)N &&\text{[\ref{hc:facts:matrix}]}\\
      &=5G+2N^TJ_PN &&\text{[\stepref{hcIncNN}]}\\
      &=5G+50J_{\cB} &&\text{[\stepref{hcIncCol}, $5\cdot5=25$].}
  \end{align*}\qed
 \end{proof}
 \step{hcIncGJ}{$GJ_{\cB}=35J_{\cB}$.}
 \begin{proof}\pf\ Steps \stepref{hcIncRow} and \stepref{hcIncCol}
  give $G\Zj_{\cB}=N^T(7\Zj_P)=35\Zj_{\cB}$;
  each column of $J$ is $\Zj$, by \ref{hc:facts:matrix}.~\qed\end{proof}
 \step{hcIncJG}{$J_{\cB}G=35J_{\cB}$.}
 \begin{proof}\pf\ Transpose \stepref{hcIncGJ}, using the symmetry
  in \stepref{hcIncG} and \ref{hc:facts:matrix}.~\qed\end{proof}
 \step{hcIncJJ}{$J_{\cB}^2=21J_{\cB}$.}
 \begin{proof}\pf\ Each entry sums $|\cB|=21$ ones, by
  Definition~\ref{hc:def:design} and \ref{hc:facts:matrix}.~\qed\end{proof}
 \qedstep
 \begin{proof}\pf\ Steps \stepref{hcIncNN}--\stepref{hcIncJJ}
  give every asserted identity.~\qed\end{proof}
\end{proof}

\begin{namedresult}{Lemma}{(intersection moments)}\label{hc:lem:moments}
 \lean{Design.intersection_moments_15_5_2}
 \uses{hc:def:design,hc:facts:finite}
 \leavevmode\par
 \assume{\begin{pfenum}
  \item $P,\cB$ are finite types with decidable equality and
  $D$ is a design with parameters $(15,5,21,7,2)$.
  \item $B\in\cB$ and $n_i\triangleq|\{B'\ne B:|B\cap B'|=i\}|$ for $0\le i\le5$.
 \end{pfenum}}
 \prove{Over $\ints$,
 $\sum_{i=0}^5n_i=20$, $\sum_{i=0}^5in_i=30$, and
 $\sum_{i=0}^5i(i-1)n_i=20$.
 In particular $\sum_{B'\ne B}|B\cap B'|=30$.}
\end{namedresult}
\begin{proof}
 \pfsketch\ We will count labels, incidences, and ordered pairs of
 distinct points. The moment statements will be cited outside this proof.
 \step{hcMomZero}{$\sum_{i=0}^5n_i=20$.}
 \begin{proof}\pf\ Intersection size is between zero and five;
  the fibres partition the $21-1$ other labels, by
  Definition~\ref{hc:def:design} and \ref{hc:facts:finite}.~\qed\end{proof}
 \step{hcMomOne}{$\sum_{i=0}^5in_i=30$ and
  $\sum_{B'\ne B}|B\cap B'|=30$.}
 \begin{proof}\pf\ Count $(p,B')$ with $p\in B\cap B'$ and $B'\ne B$.
  For each of five points, six of its seven labels remain.
  Grouping by $B'$ and then by $i$ gives both sums, using
  Definition~\ref{hc:def:design} and \ref{hc:facts:finite}.~\qed\end{proof}
 \step{hcMomTwo}{$\sum_{i=0}^5i(i-1)n_i=20$.}
 \begin{proof}\pf\ Count $(p,q,B')$ with $p,q\in B\cap B'$, $p\ne q$,
  $B'\ne B$. Each of $5\cdot4$ ordered pairs in $B$ has exactly
  one further label; a label with intersection size $i$ contributes
  $i(i-1)$, using Definition~\ref{hc:def:design} and
  \ref{hc:facts:finite}.~\qed\end{proof}
 \qedstep
 \begin{proof}\pf\ Steps \stepref{hcMomZero}, \stepref{hcMomOne},
  and \stepref{hcMomTwo} prove the four assertions.~\qed\end{proof}
\end{proof}

\begin{namedresult}{Lemma}{(block intersections)}\label{hc:lem:intersections}
 \textup{[\cite{HC}, Lemma 4.1 at $k=5$.]}
 \lean{Design.inter_card_eq_one_or_two}\lean{Design.card_neighbors}
 \lean{Design.card_filter_inter_card_eq_two}
 \uses{hc:def:design,hc:def:graph,hc:lem:moments,hc:facts:finite}
 \leavevmode\par
 \assume{\begin{pfenum}
  \item $P,\cB$ are finite types with decidable equality.
  \item $D$ is a design with parameters $(15,5,21,7,2)$.
 \end{pfenum}}
 \prove{Distinct labels meet in one or two points. Every block has
 exactly ten other labels of each kind; in particular $|\Gamma(B)|=10$.}
\end{namedresult}
\begin{proof}
 \pfsketch\ A nonnegative linear combination of the moments excludes
 intersection sizes other than one and two.
 \step{hcInterBlock}{\assume{$B\in\cB$, with $n_i$ as in
  Lemma~\ref{hc:lem:moments}.}
  \prove{Every other label meets $B$ in one or two points,
  and exactly ten labels have each intersection size.}}
 \begin{proof}
  \step{hcInterPolynomial}{$\sum_{i=0}^5(i-1)(i-2)n_i=0$.}
  \begin{proof}\pf\ The polynomial equals $i(i-1)-2i+2$;
   Lemma~\ref{hc:lem:moments} gives $20-2\cdot30+2\cdot20=0$.~\qed\end{proof}
  \step{hcInterSupport}{$n_0=n_3=n_4=n_5=0$.}
  \begin{proof}
   \step{hcInterCoeffs}{The coefficients $(i-1)(i-2)$, in order
    $i=0,1,2,3,4,5$, are $2,0,0,2,6,12$.}
   \begin{proof}\pf\ Direct integer substitution in the polynomial
    of \stepref{hcInterPolynomial}.~\qed\end{proof}
   \qedstep
   \begin{proof}\pf\ The $n_i$ are nonnegative cardinalities.
    Apply the zero-sum criterion of \ref{hc:facts:finite} to
    \stepref{hcInterPolynomial} and \stepref{hcInterCoeffs}; the
    four indicated coefficients are positive.~\qed\end{proof}
  \end{proof}
  \step{hcInterCounts}{$n_1=n_2=10$.}
  \begin{proof}\pf\ Lemma~\ref{hc:lem:moments} and
   \stepref{hcInterSupport} give $n_1+n_2=20$ and $n_1+2n_2=30$.
   Subtract the equations.~\qed\end{proof}
  \qedstep
  \begin{proof}\pf\ The zero fibres of \stepref{hcInterSupport}
   exclude the other intersection sizes; \stepref{hcInterCounts}
   and Definition~\ref{hc:def:graph} give the multiplicities and degree.~\qed\end{proof}
 \end{proof}
 \qedstep
 \begin{proof}\pf\ Step \stepref{hcInterBlock} applies to every $B$;
  Definition~\ref{hc:def:graph} identifies its ten one-point
  intersections with $\Gamma(B)$.~\qed\end{proof}
\end{proof}

\begin{namedresult}{Lemma}{(common-neighbour counts)}\label{hc:lem:common}
 \textup{[\cite{HC}, Lemma 4.2 at $k=5$.]}
 \lean{Design.intersection_pair_counts}\lean{Design.card_commonNeighbors}
 \uses{hc:def:design,hc:def:matrices,hc:def:graph,hc:lem:incidence,
 hc:lem:moments,hc:lem:intersections,hc:facts:finite,hc:facts:matrix}
 \leavevmode\par
 \assume{\begin{pfenum}
  \item $P,\cB$ are finite types with decidable equality and
  $D$ is a design with parameters $(15,5,21,7,2)$.
  \item $B_1\ne B_2$ are block labels, $s\triangleq|B_1\cap B_2|$, and
  $x_{ij}\triangleq|\{B\notin\{B_1,B_2\}:|B_1\cap B|=i,\ |B_2\cap B|=j\}|$.
 \end{pfenum}}
 \prove{The four counts are
 \[
 \begin{array}{c|rrrr}
 s&x_{11}&x_{12}&x_{21}&x_{22}\\\hline
 1&5&4&4&6\\2&4&6&6&3
 \end{array}.
 \]
 Thus adjacent vertices have five common neighbours and distinct
 nonadjacent vertices have four.}
\end{namedresult}
\begin{proof}
 \pfsketch\ We will expose the four filtered-sum equations, including
 both terms removed from the entry of $G^2$, before solving them.
 \step{hcCommonTotal}{The four fibres partition
  $\cB\diff\{B_1,B_2\}$, and $x_{11}+x_{12}+x_{21}+x_{22}=19$.}
 \begin{proof}\pf\ Lemma~\ref{hc:lem:intersections} partitions the
  $21-2$ remaining labels into the four disjoint fibres; use
  \asm[2]{0}, \ref{hc:facts:finite}, and
  Definition~\ref{hc:def:design}.~\qed\end{proof}
 \step{hcCommonFirst}{$x_{11}+x_{12}+2x_{21}+2x_{22}=30-s$.}
 \begin{proof}\pf\ From the last conclusion of Lemma~\ref{hc:lem:moments}
  for $B_1$, remove the term $B_2$ of size $s$. The four fibres have
  first weights $1,1,2,2$, using \stepref{hcCommonTotal},
  Lemma~\ref{hc:lem:intersections}, \asm[2]{0}, and \ref{hc:facts:finite}.~\qed\end{proof}
 \step{hcCommonSecond}{$x_{11}+2x_{12}+x_{21}+2x_{22}=30-s$.}
 \begin{proof}\pf\ Apply Lemma~\ref{hc:lem:moments} to $B_2$ and remove
  the term $B_1$ of size $s$. The second weights are $1,2,1,2$,
  using \stepref{hcCommonTotal}, Lemma~\ref{hc:lem:intersections},
  \asm[2]{0}, and \ref{hc:facts:finite}.~\qed\end{proof}
 \step{hcCommonProduct}{$x_{11}+2x_{12}+2x_{21}+4x_{22}=50-5s$.}
 \begin{proof}
  \step{hcCommonEntry}{$(G^2)(B_1,B_2)=5s+50$.}
  \begin{proof}\pf\ Substitute $G(B_1,B_2)=s$ into the identity
   $G^2=5G+50J$ of Lemma~\ref{hc:lem:incidence}.~\qed\end{proof}
  \step{hcCommonRemoved}{The terms indexed by $B_1$ and $B_2$ in
   $(G^2)(B_1,B_2)$ are each $5s$.}
  \begin{proof}\pf\ The products are $G(B_1,B_1)G(B_1,B_2)$ and
   $G(B_1,B_2)G(B_2,B_2)$; both diagonals equal five by
   Definition~\ref{hc:def:design} and Lemma~\ref{hc:lem:incidence}.~\qed\end{proof}
  \qedstep
  \begin{proof}\pf\ Expand the matrix entry using \ref{hc:facts:matrix},
   remove the two distinct terms of \stepref{hcCommonRemoved} from
   \stepref{hcCommonEntry}, then group the remaining terms by their
   weights $ij$ using \stepref{hcCommonTotal},
   Lemma~\ref{hc:lem:intersections}, \asm[2]{0}, and
   \ref{hc:facts:finite}. This gives
   $(5s+50)-10s=50-5s$.~\qed\end{proof}
 \end{proof}
 \step{hcCommonSolve}{$x_{22}=9-3s$, $x_{12}=x_{21}=2+2s$, and $x_{11}=6-s$.}
 \begin{proof}\pf\ Subtract \stepref{hcCommonFirst} and
  \stepref{hcCommonSecond} from \stepref{hcCommonProduct}, then add
  \stepref{hcCommonTotal}, to get $x_{22}=9-3s$.
  Subtract \stepref{hcCommonTotal} and this value from each weighted
  marginal, then use the total for $x_{11}$.~\qed\end{proof}
 \step{hcCommonOne}{\case{$s=1$.} The count vector is $(5,4,4,6)$.}
 \begin{proof}\pf\ Substitute \asm{1} in \stepref{hcCommonSolve}.~\qed\end{proof}
 \step{hcCommonTwo}{\case{$s=2$.} The count vector is $(4,6,6,3)$.}
 \begin{proof}\pf\ Substitute \asm{1} in \stepref{hcCommonSolve}.~\qed\end{proof}
 \step{hcCommonNeighbours}{$\Gamma(B_1)\cap\Gamma(B_2)$ is exactly
  the $x_{11}$ fibre; $B_1\sim B_2$ if and only if $s=1$.}
 \begin{proof}\pf\ By Definition~\ref{hc:def:graph}, membership
  in both neighbourhoods is intersection size one with each block;
  irreflexivity excludes $B_1,B_2$. Conversely the fibre excludes
  these two labels and gives both adjacencies. For the pair itself,
  \asm[2]{0} and Lemma~\ref{hc:lem:intersections} leave precisely
  $s=1$ for adjacency and $s=2$ for nonadjacency.~\qed\end{proof}
 \qedstep
 \begin{proof}\pf\ Lemma~\ref{hc:lem:intersections} makes the cases
  \stepref{hcCommonOne}, \stepref{hcCommonTwo} exhaustive.
  Step \stepref{hcCommonNeighbours} identifies their $x_{11}$
  counts with the asserted common-neighbour counts.~\qed\end{proof}
\end{proof}

\subsubsection{Connor's matrix and its rank}

\begin{namedresult}{Lemma}{(a scalar matrix plus a constant matrix)}\label{hc:lem:det}
 \lean{Matrix.det_smul_one_add_smul_J}
 \uses{hc:facts:det}
 \leavevmode\par
 \assume{\begin{pfenum}
  \item $U$ is a finite type with decidable equality, $n=|U|\ge1$.
  \item $\alpha,\beta\in\rats$ and $\alpha\ne0$.
 \end{pfenum}}
 \prove{$\det(\alpha I_U+\beta J_U)=\alpha^{n-1}(\alpha+n\beta)$.}
\end{namedresult}
\begin{proof}
 \pfsketch\ Factor out the nonzero scalar and apply the matrix
 determinant lemma from \ref{hc:facts:det} to a rank-one product.
 \step{hcDetFactor}{$\alpha I+\beta J=
  \alpha\bigl(I+(\tfrac\beta\alpha\Zj_U)\Zj_U^T\bigr)$.}
 \begin{proof}\pf\ Every entry of $\Zj_U\Zj_U^T$ is one.
  Distribute $\alpha$ and cancel $\alpha\ne0$ from \asm[2]{0}.~\qed\end{proof}
 \step{hcDetRankOne}{\new\ $\det\bigl(I+(\tfrac\beta\alpha\Zj_U)\Zj_U^T\bigr)
  =1+n\tfrac\beta\alpha$.}
 \begin{proof}\pf\ Apply
  \path{Matrix.det_one_add_replicateCol_mul_replicateRow}
  from \ref{hc:facts:det}, with singleton index \texttt{Unit},
  $\Zu=\tfrac\beta\alpha\Zj_U$, and $\Zv=\Zj_U$.
  The dot product sums $n$ copies of $\beta/\alpha$.~\qed\end{proof}
 \step{hcDetScale}{$\det(\alpha I+\beta J)=\alpha^{n-1}(\alpha+n\beta)$.}
 \begin{proof}\pf\ By \path{Matrix.det_smul} in \ref{hc:facts:det},
  steps \stepref{hcDetFactor} and \stepref{hcDetRankOne} give
  $\alpha^n(1+n\beta/\alpha)$.
  Use $n=(n-1)+1$ from \asm[1]{0} and cancel $\alpha\ne0$.~\qed\end{proof}
 \qedstep
 \begin{proof}\pf\ Step \stepref{hcDetScale}.~\qed\end{proof}
\end{proof}

\begin{namedresult}{Lemma}{(Connor identities and entries)}\label{hc:lem:Connor}
 \lean{Design.C_transpose}\lean{Design.C_apply}
 \lean{Design.C_mul_C}\lean{Design.C_mul_incMatrix_transpose}
 \uses{hc:def:design,hc:def:matrices,hc:def:graph,hc:lem:incidence,
 hc:lem:intersections,hc:facts:matrix}
 \leavevmode\par
 \assume{\begin{pfenum}
  \item $P,\cB$ are finite types with decidable equality.
  \item $D$ is a design with parameters $(15,5,21,7,2)$.
 \end{pfenum}}
 \prove{Over $\rats$, $C^T=C$, $C^2=35C$, $CN^T=0$, and
 \[
 C(B,B')=\begin{cases}
 10,&B=B',\\3,&B\sim B',\\-4,&B\ne B'\text{ and }B\not\sim B'.
 \end{cases}
 \]}
\end{namedresult}
\begin{proof}
 \pfsketch\ We will substitute the identities for $G$ into
 $C=35I-7G+10J$ and check its three entry types separately.
 \step{hcCSym}{$C^T=C$.}
 \begin{proof}\pf\ $I,J,G$ are symmetric by
  Lemma~\ref{hc:lem:incidence}; transpose the defining linear
  combination in Definition~\ref{hc:def:matrices}.~\qed\end{proof}
 \step{hcCDiag}{\assume{$B=B'$.}\prove{$C(B,B')=10$.}}
 \begin{proof}\pf\ Lemma~\ref{hc:lem:incidence} and
  Definition~\ref{hc:def:design} give $G(B,B)=5$, so
  Definition~\ref{hc:def:matrices} gives $35-7\cdot5+10=10$.~\qed\end{proof}
 \step{hcCAdj}{\assume{$B\sim B'$.}\prove{$C(B,B')=3$.}}
 \begin{proof}\pf\ Definition~\ref{hc:def:graph} gives distinct
  labels and intersection size one. Definition~\ref{hc:def:matrices}
  and Lemma~\ref{hc:lem:incidence} give $0-7\cdot1+10=3$.~\qed\end{proof}
 \step{hcCNonadj}{\assume{$B\ne B'$ and $B\not\sim B'$.}
  \prove{$C(B,B')=-4$.}}
 \begin{proof}\pf\ Lemma~\ref{hc:lem:intersections} and
  Definition~\ref{hc:def:graph} force intersection size two.
  Definition~\ref{hc:def:matrices} gives $0-7\cdot2+10=-4$.~\qed\end{proof}
 \step{hcCSquare}{$C^2=35C$.}
 \begin{proof}
  \pf\ In the following calculation $I,G,J$ are indexed by $\cB$:
  \begin{align*}
   C^2&=1225I+49G^2+100J^2-490G+700J-140GJ
     &&\text{[\ref{hc:lem:incidence}, $GJ=JG$]}\\
   &=1225I+49(5G+50J)-490G-2100J
     &&\text{[\ref{hc:lem:incidence}, arithmetic]}\\
   &=1225I-245G+350J &&\text{[ring arithmetic]}\\
   &=35C &&\text{[\ref{hc:def:matrices}].}
  \end{align*}\qed
 \end{proof}
 \step{hcCZero}{$CN^T=0$.}
 \begin{proof}
  \step{hcCOnes}{$N^TJ_P=5J_{\cB,P}$ and $J_{\cB}N^T=7J_{\cB,P}$.}
  \begin{proof}\pf\ Each entry in the first product is a column
   sum of $N$, and each in the second a row sum; apply
   Lemma~\ref{hc:lem:incidence} and \ref{hc:facts:matrix}.~\qed\end{proof}
  \qedstep
  \begin{proof}
   \pf\ By associativity from \ref{hc:facts:matrix},
   \begin{align*}
    CN^T&=35N^T-7N^T(NN^T)+10J_{\cB}N^T
      &&\text{[\ref{hc:def:matrices}]}\\
    &=35N^T-7N^T(5I_P+2J_P)+70J_{\cB,P}
      &&\text{[\ref{hc:lem:incidence}, \stepref{hcCOnes}]}\\
    &=0 &&\text{[\stepref{hcCOnes}, ring arithmetic].}
   \end{align*}\qed
  \end{proof}
 \end{proof}
 \qedstep
 \begin{proof}\pf\ Steps \stepref{hcCDiag}, \stepref{hcCAdj},
  \stepref{hcCNonadj} cover every pair of indices.
  Steps \stepref{hcCSym}, \stepref{hcCSquare}, \stepref{hcCZero}
  give the other conclusions.~\qed\end{proof}
\end{proof}

\begin{namedresult}{Lemma}{(quadratic inequality and rank bound)}\label{hc:lem:rank}
 \lean{Design.C_quadratic_nonneg}\lean{Design.rank_C_le_six}
 \uses{hc:def:design,hc:def:matrices,hc:lem:incidence,hc:lem:Connor,
 hc:lem:det,hc:facts:matrix,hc:facts:rank}
 \leavevmode\par
 \assume{\begin{pfenum}
  \item $P,\cB$ are finite types with decidable equality.
  \item $D$ is a design with parameters $(15,5,21,7,2)$.
 \end{pfenum}}
 \prove{For every $\Zx\in\rats^{\cB}$,
 $35\Zx^TC\Zx=(C\Zx)^T(C\Zx)\ge0$, so $\Zx^TC\Zx\ge0$;
 also $\operatorname{rk}C\le6$.}
\end{namedresult}
\begin{proof}
 \pfsketch\ The equation $C^2=35C$ gives a sum of squares. The
 equation $CN^T=0$ gives the rank bound after proving $N$ has rank fifteen.
 \step{hcRankQuad}{For every $\Zx$, $35\Zx^TC\Zx=(C\Zx)^T(C\Zx)\ge0$.}
 \begin{proof}\pf\ Lemma~\ref{hc:lem:Connor} gives $35C=C^TC$.
  Multiply by $\Zx^T,\Zx$, using \ref{hc:facts:matrix}; the resulting
  value is a sum of rational squares.~\qed\end{proof}
 \step{hcRankNonneg}{For every $\Zx$, $\Zx^TC\Zx\ge0$.}
 \begin{proof}\pf\ Divide \stepref{hcRankQuad} by $35>0$.~\qed\end{proof}
 \step{hcRankDet}{$\det(NN^T)=5^{14}\cdot35\ne0$.}
 \begin{proof}\pf\ First cast Lemma~\ref{hc:lem:incidence} to $\rats$
  by \ref{hc:facts:matrix}. Apply Lemma~\ref{hc:lem:det}.
  Definition~\ref{hc:def:design} gives
  $|P|=15$; take $\alpha=5\ne0$ and $\beta=2$.~\qed\end{proof}
 \step{hcRankN}{$\operatorname{rk}N=\operatorname{rk}N^T=15$.}
 \begin{proof}
  \step{hcRankFull}{$\operatorname{rk}(NN^T)=15$.}
  \begin{proof}\pf\ The determinant in \stepref{hcRankDet} is a
   unit in $\rats$. The unit determinant and full-rank criteria of
   \ref{hc:facts:rank} give rank $|P|=15$.~\qed\end{proof}
  \qedstep
  \begin{proof}\pf\ Apply the two transpose-rank identities of
   \ref{hc:facts:rank} over $\rats$ to \stepref{hcRankFull}.~\qed\end{proof}
 \end{proof}
 \step{hcRankBound}{$\operatorname{rk}C\le6$.}
 \begin{proof}\pf\ By $CN^T=0$ in Lemma~\ref{hc:lem:Connor} and
  the zero-product inequality of \ref{hc:facts:rank},
  $\operatorname{rk}C+\operatorname{rk}N^T\le|\cB|=21$.
  Substitute \stepref{hcRankN} and subtract fifteen.~\qed\end{proof}
 \qedstep
 \begin{proof}\pf\ Steps \stepref{hcRankQuad}, \stepref{hcRankNonneg},
  and \stepref{hcRankBound}.~\qed\end{proof}
\end{proof}

\begin{namedresult}{Lemma}{(principal matrices on labelled subsets)}\label{hc:lem:principal}
 \lean{Design.C_submatrix_quadratic_nonneg}
 \lean{Design.rank_C_submatrix_le_six}\lean{Design.det_C_submatrix_eq_zero}
 \uses{hc:def:design,hc:def:matrices,hc:lem:rank,hc:facts:finite,
 hc:facts:matrix,hc:facts:rank}
 \leavevmode\par
 \assume{\begin{pfenum}
  \item $P,\cB,U$ are finite types with decidable equality and
  $D$ is a design with parameters $(15,5,21,7,2)$.
  \item $e:U\to\cB$ is injective, and $C_e(u,w)\triangleq C(e(u),e(w))$.
 \end{pfenum}}
 \prove{For all $\Zx\in\rats^U$, $\Zx^TC_e\Zx\ge0$;
 $\operatorname{rk}C_e\le6$; if $|U|\ge7$, then $\det C_e=0$.}
\end{namedresult}
\begin{proof}
 \pfsketch\ We will extend vectors by zero along the injection.
 Rank monotonicity also applies to the specified submatrix, so later
 enumerations need no unstated transport of singularity.
 \step{hcPrinExtension}{For $\Zx\in\rats^U$, define
  $\widehat\Zx(e(u))=\Zx(u)$ and $\widehat\Zx(B)=0$ off the image.
  This defines $\widehat\Zx\in\rats^{\cB}$.}
 \begin{proof}\pf\ Injectivity in \asm[2]{0} gives at most one
  preimage; the two domains partition $\cB$.~\qed\end{proof}
 \step{hcPrinValue}{$\widehat\Zx^TC\widehat\Zx=\Zx^TC_e\Zx\ge0$.}
 \begin{proof}\pf\ Expand the quadratic sum using \ref{hc:facts:matrix}.
  Terms off the image vanish by \stepref{hcPrinExtension}; change both
  finite indices along $e$ using \ref{hc:facts:finite}. Apply
  Lemma~\ref{hc:lem:rank} to $\widehat\Zx$.~\qed\end{proof}
 \step{hcPrinRank}{$\operatorname{rk}C_e\le6$.}
 \begin{proof}\pf\ Rank monotonicity for submatrices in
  \ref{hc:facts:rank} bounds this rank by $\operatorname{rk}C$;
  use Lemma~\ref{hc:lem:rank}.~\qed\end{proof}
 \step{hcPrinSingular}{\assume{$|U|\ge7$.}\prove{$\det C_e=0$.}}
 \begin{proof}\pf\ If the determinant were nonzero, the full-rank
  criterion in \ref{hc:facts:rank} would give
  $\operatorname{rk}C_e=|U|\ge7$, contradicting \stepref{hcPrinRank}.~\qed\end{proof}
 \qedstep
 \begin{proof}\pf\ Steps \stepref{hcPrinValue}, \stepref{hcPrinRank},
  and \stepref{hcPrinSingular}.~\qed\end{proof}
\end{proof}

\subsubsection{Finite certificates and local restrictions}

\begin{namedresult}{Definition}{(the C19 matrices)}\label{hc:def:certificates}
 \lean{Computations.C19.clawMatrix}\lean{Computations.C19.fiveBlockMatrix}
 \lean{Computations.C19.Q1}\lean{Computations.C19.Q2}
 \lean{Computations.C19.K1}\lean{Computations.C19.K2}\uses{}
 The following matrices have integer entries; their rational casts
 are denoted by the same letters. Indices in the displays run from
 $1$ to the indicated order, corresponding to Lean's
 $\mathrm{Fin}\,n$ in increasing order.

\[
 Q_{\rm claw}=
% BEGIN MATRIX claw
 \begin{pmatrix}
 10&3&3&3\\3&10&-4&-4\\3&-4&10&-4\\3&-4&-4&10
 \end{pmatrix}
% END MATRIX claw
\]

\[
 Q_{\rm five}=
% BEGIN MATRIX five
 \begin{pmatrix}
 10&3&3&3&-4\\3&10&3&-4&3\\3&3&10&-4&3\\
 3&-4&-4&10&3\\-4&3&3&3&10
 \end{pmatrix}
% END MATRIX five
\]

\[
 Q_1=
% BEGIN MATRIX Q1
 \begin{pmatrix}
 10&3&3&3&-4&-4&-4\\
 3&10&3&-4&3&3&3\\
 3&3&10&-4&3&3&3\\
 3&-4&-4&10&-4&-4&-4\\
 -4&3&3&-4&10&3&3\\
 -4&3&3&-4&3&10&3\\
 -4&3&3&-4&3&3&10
 \end{pmatrix}
% END MATRIX Q1
\]

\[
 Q_2=
% BEGIN MATRIX Q2
 \begin{pmatrix}
 10&3&3&3&3&-4&-4\\
 3&10&3&-4&-4&3&3\\
 3&3&10&-4&-4&3&3\\
 3&-4&-4&10&3&-4&-4\\
 3&-4&-4&3&10&-4&-4\\
 -4&3&3&-4&-4&10&3\\
 -4&3&3&-4&-4&3&10
 \end{pmatrix}
% END MATRIX Q2
\]

\[
 K_1=
% BEGIN MATRIX K1
 \begin{pmatrix}
 20&-12&-12&-6&8&8&8\\
 -12&9&7&4&-5&-5&-5\\
 -12&7&9&4&-5&-5&-5\\
 -6&4&4&4&-2&-2&-2\\
 8&-5&-5&-2&5&3&3\\
 8&-5&-5&-2&3&5&3\\
 8&-5&-5&-2&3&3&5
 \end{pmatrix}
% END MATRIX K1
\]

\[
 K_2=
% BEGIN MATRIX K2
 \begin{pmatrix}
 17&-11&-11&-7&-7&6&6\\
 -11&9&7&5&5&-4&-4\\
 -11&7&9&5&5&-4&-4\\
 -7&5&5&5&3&-2&-2\\
 -7&5&5&3&5&-2&-2\\
 6&-4&-4&-2&-2&4&2\\
 6&-4&-4&-2&-2&2&4
 \end{pmatrix}
% END MATRIX K2
\]
\end{namedresult}

\begin{namedresult}{Computation}{(C19: quadratic values and inverse certificates)}\label{hc:comp:C19}
 \lean{Computations.C19.claw_quadratic}
 \lean{Computations.C19.five_block_quadratic}
 \lean{Computations.C19.left_inverse_Q1}\lean{Computations.C19.left_inverse_Q2}
 \uses{hc:def:certificates,hc:facts:matrix}
 \leavevmode\par
 \pflet{$\Zx\triangleq(-1,1,1,1)^T$ and
 $\Zy\triangleq(-1,1,1,1,-1)^T$.}
 \prove{$\Zx^TQ_{\rm claw}\Zx=-2$, $\Zy^TQ_{\rm five}\Zy=-4$, and
 $K_jQ_j=14I_7$ for each $j\in\{1,2\}$.}
\end{namedresult}
\begin{proof}
 \pfsketch\ These are finite integer certificates, suitable for
 kernel-checked arithmetic. The script \path{blueprint/checks/hallconnor/check_proof.py} reads every matrix
 from Definition~\ref{hc:def:certificates} and recomputes every entry.
 \step{hcCertClaw}{$\Zx^TQ_{\rm claw}\Zx=-2$.}
 \begin{proof}\pf\ From Definition~\ref{hc:def:certificates}, the
  diagonal sum is $40$ and the upper-triangle sum is $-9-12=-21$.
  The quadratic sum is $40+2(-21)=-2$, by \ref{hc:facts:matrix}.~\qed\end{proof}
 \step{hcCertFive}{$\Zy^TQ_{\rm five}\Zy=-4$.}
 \begin{proof}\pf\ From Definition~\ref{hc:def:certificates}, the
  diagonal sum is $50$. The upper-triangle row contributions are
  $-13,-4,-7,-3$, totalling $-27$. Thus the value is
  $50+2(-27)=-4$, by \ref{hc:facts:matrix}.~\qed\end{proof}
 \step{hcCertOne}{$K_1Q_1=14I_7$.}
 \begin{proof}\pf\ For all $u,w\in\{1,\ldots,7\}$, substitute the
  integer entries in Definition~\ref{hc:def:certificates} into
  $\sum_{i=1}^7K_1(u,i)Q_1(i,w)$, using \ref{hc:facts:matrix}.
  The seven diagonal sums are $14$ and all $42$ other sums are zero.
  For example the first row starts with
  $200-36-36-18-32-32-32=14$ and
  $60-120-36+24+24+24+24=0$.~\qed\end{proof}
 \step{hcCertTwo}{$K_2Q_2=14I_7$.}
 \begin{proof}\pf\ For all $u,w\in\{1,\ldots,7\}$, the sums
  $\sum_{i=1}^7K_2(u,i)Q_2(i,w)$ from
  Definition~\ref{hc:def:certificates} are $14$ on the diagonal
  and zero off it, by integer multiplication and \ref{hc:facts:matrix}.
  The first two sums are
  $170-33-33-21-21-24-24=14$ and
  $51-110-33+28+28+18+18=0$.~\qed\end{proof}
 \qedstep
 \begin{proof}\pf\ Steps \stepref{hcCertClaw}, \stepref{hcCertFive},
  \stepref{hcCertOne}, and \stepref{hcCertTwo}.~\qed\end{proof}
\end{proof}

\begin{namedresult}{Lemma}{(the two nonsingular matrices)}\label{hc:lem:inverse}
 \lean{Computations.C19.det_Q1_ne_zero}\lean{Computations.C19.det_Q2_ne_zero}
 \uses{hc:def:certificates,hc:comp:C19,hc:facts:matrix,hc:facts:det}
 \leavevmode\par
 \assume{$j\in\{1,2\}$, with $Q_j,K_j$ as in Definition~\ref{hc:def:certificates}.}
 \prove{$\det Q_j\ne0$ over $\rats$.}
\end{namedresult}
\begin{proof}
 \pfsketch\ The certified product has nonzero determinant.
 \step{hcInvProduct}{\new\ $(\det K_j)(\det Q_j)=14^7$.}
 \begin{proof}\pf\ Cast the integral identity of
  Computation~\ref{hc:comp:C19} to $\rats$ by \ref{hc:facts:matrix},
  then take determinants;
  multiplicativity and scalar multiplication in \ref{hc:facts:det}
  give $\det(14I_7)=14^7$. This replaces the ordinary proof's
  zero-kernel argument with the determinant of the same certificate.~\qed\end{proof}
 \qedstep
 \begin{proof}\pf\ If $\det Q_j=0$, the left side of
  \stepref{hcInvProduct} vanishes, whereas $14^7\ne0$ in $\rats$.~\qed\end{proof}
\end{proof}

\begin{namedresult}{Lemma}{(no induced claw)}\label{hc:lem:claw}
 \lean{Design.Z_clawFree}
 \uses{hc:def:design,hc:def:matrices,hc:def:graph,hc:def:certificates,hc:lem:Connor,
 hc:lem:principal,hc:comp:C19}
 \leavevmode\par
 \assume{\begin{pfenum}
  \item $P,\cB$ are finite types with decidable equality.
  \item $D$ is a design with parameters $(15,5,21,7,2)$.
 \end{pfenum}}
 \prove{There are no four distinct blocks $B,X,Y,W$ for which
 $B\sim X$, $B\sim Y$, $B\sim W$, and $X,Y,W$ are pairwise nonadjacent.}
\end{namedresult}
\begin{proof}
 \pfsketch\ The four specified labels would give a negative
 quadratic value of a principal matrix of $C$.
 \step{hcClawBad}{\assume{Four blocks $B,X,Y,W$ satisfy the forbidden pattern.}
  \prove{\textsc{False}.}}
 \begin{proof}
  \step{hcClawIndex}{The ordered list $(B,X,Y,W)$ defines an
   injection $e:\mathrm{Fin}\,4\to\cB$.}
  \begin{proof}\pf\ The blocks are distinct by \asm{1}.~\qed\end{proof}
  \step{hcClawDiag}{$C_e(i,i)=10$ for each $i$.}
  \begin{proof}\pf\ The diagonal clause of Lemma~\ref{hc:lem:Connor}.~\qed\end{proof}
  \step{hcClawCentre}{All entries from the centre to a leaf are $3$.}
  \begin{proof}\pf\ The three adjacencies in \asm{1} and the
   adjacency clause of Lemma~\ref{hc:lem:Connor}.~\qed\end{proof}
  \step{hcClawLeaves}{All entries between distinct leaves are $-4$.}
  \begin{proof}\pf\ Their distinctness and nonadjacency in \asm{1}
   give the last clause of Lemma~\ref{hc:lem:Connor}.~\qed\end{proof}
  \step{hcClawMatrix}{$C_e=Q_{\rm claw}$.}
  \begin{proof}\pf\ The diagonal, centre-leaf, and leaf-leaf entries
   are exactly those of Definition~\ref{hc:def:certificates}, by
   \stepref{hcClawDiag}, \stepref{hcClawCentre}, \stepref{hcClawLeaves}.
   Symmetry in Lemma~\ref{hc:lem:Connor} supplies the reversed entries.~\qed\end{proof}
  \qedstep
  \begin{proof}\pf\ Apply Lemma~\ref{hc:lem:principal} with
   \stepref{hcClawIndex} to the vector of Computation~\ref{hc:comp:C19}.
   Step \stepref{hcClawMatrix} makes its value $-2\ge0$, a contradiction.~\qed\end{proof}
 \end{proof}
 \qedstep
 \begin{proof}\pf\ Step \stepref{hcClawBad} negates the forbidden existence.~\qed\end{proof}
\end{proof}

\begin{namedresult}{Lemma}{(the five-block restriction)}\label{hc:lem:five}
 \textup{[\cite{HC}, Lemma 4.5 at $k=5$.]}
 \lean{Design.not_adj_of_five_block_pattern}
 \uses{hc:def:design,hc:def:matrices,hc:def:graph,hc:def:certificates,hc:lem:Connor,
 hc:lem:principal,hc:comp:C19}
 \leavevmode\par
 \assume{\begin{pfenum}
  \item $P,\cB$ are finite types with decidable equality and
  $D$ is a design with parameters $(15,5,21,7,2)$.
  \item $B_1,B_2,B_3,Y,W$ are five distinct blocks; $B_1,B_2,B_3$ form a triangle.
  \item $Y\sim B_1$, $Y\not\sim B_2,B_3$, and
  $W\sim B_2,B_3$, $W\not\sim B_1$.
 \end{pfenum}}
 \prove{$Y\not\sim W$.}
\end{namedresult}
\begin{proof}
 \pfsketch\ An edge $YW$ gives the second negative quadratic certificate.
 \step{hcFiveBad}{\assume{$Y\sim W$.}\prove{\textsc{False}.}}
 \begin{proof}
  \step{hcFiveIndex}{The ordered list $(B_1,B_2,B_3,Y,W)$ defines an
   injection $e:\mathrm{Fin}\,5\to\cB$.}
  \begin{proof}\pf\ Distinctness is \asm[2]{0}.~\qed\end{proof}
  \step{hcFiveDiag}{Every diagonal entry of $C_e$ is $10$.}
  \begin{proof}\pf\ Lemma~\ref{hc:lem:Connor}.~\qed\end{proof}
  \step{hcFiveTriangle}{The three entries between $B_1,B_2,B_3$ are $3$.}
  \begin{proof}\pf\ The triangle in \asm[2]{0} and
   Lemma~\ref{hc:lem:Connor}.~\qed\end{proof}
  \step{hcFiveY}{The entries from $Y$ to $(B_1,B_2,B_3)$ are $(3,-4,-4)$.}
  \begin{proof}\pf\ Apply the adjacency and nonadjacency clauses of
   Lemma~\ref{hc:lem:Connor} to the three stipulated relations in
   \asm[3]{0}; their labels are distinct by \asm[2]{0}.~\qed\end{proof}
  \step{hcFiveW}{The entries from $W$ to $(B_1,B_2,B_3)$ are $(-4,3,3)$.}
  \begin{proof}\pf\ Apply Lemma~\ref{hc:lem:Connor} to the three
   stipulated relations of $W$ in \asm[3]{0}, using \asm[2]{0}.~\qed\end{proof}
  \step{hcFiveYW}{$C(Y,W)=3$.}
  \begin{proof}\pf\ By \asm{1} and Lemma~\ref{hc:lem:Connor}.~\qed\end{proof}
  \step{hcFiveMatrix}{$C_e=Q_{\rm five}$.}
  \begin{proof}\pf\ Steps \stepref{hcFiveDiag}--\stepref{hcFiveYW}
   give every diagonal and upper-triangle entry in
   Definition~\ref{hc:def:certificates}; symmetry from
   Lemma~\ref{hc:lem:Connor} supplies the lower triangle.~\qed\end{proof}
  \qedstep
  \begin{proof}\pf\ Lemma~\ref{hc:lem:principal},
   \stepref{hcFiveIndex}, \stepref{hcFiveMatrix}, and
   Computation~\ref{hc:comp:C19} give $-4\ge0$, a contradiction.~\qed\end{proof}
 \end{proof}
 \qedstep
 \begin{proof}\pf\ Step \stepref{hcFiveBad} proves the negation.~\qed\end{proof}
\end{proof}

\subsubsection{The local six-cliques}

\begin{namedresult}{Lemma}{(splitting a neighbourhood)}\label{hc:lem:split}
 \lean{Design.exists_sixClique_neighborhood_split}
 \uses{hc:def:design,hc:def:matrices,hc:def:graph,hc:def:certificates,hc:lem:intersections,
 hc:lem:common,hc:lem:claw,hc:lem:five,hc:lem:Connor,hc:lem:principal,
 hc:lem:inverse,hc:facts:finite}
 \leavevmode\par
 \assume{\begin{pfenum}
  \item $P,\cB$ are finite types with decidable equality and
  $D$ is a design with parameters $(15,5,21,7,2)$.
  \item $B_1,B_2\in\cB$ and $B_1\sim B_2$.
 \end{pfenum}}
 \prove{There are $S_1,S_2\in\cS$ such that
 \[
 B_1,B_2\in S_1,\quad S_1\cap S_2=\{B_1\},\quad
 S_1\cup S_2=\{B_1\}\cup\Gamma(B_1).
 \]
 Moreover every $K\in\cS$ containing $B_1,B_2$ equals $S_1$.}
\end{namedresult}
\begin{proof}
 \pfsketch\ We will partition the other blocks by adjacency to $B_1,B_2$.
 The claw restriction gives two cliques and an exact edge count.
 A suitable third vertex then has a defect $j\le2$. The two positive
 defects give the certified nonsingular seven-block matrices, in
 conflict with the rank bound. Defect zero gives the split and uniqueness.
 \pflet{$\cM\triangleq\Gamma(B_1)\cap\Gamma(B_2)$,
 $\cU\triangleq\Gamma(B_1)\diff(\Gamma(B_2)\cup\{B_2\})$,
 $\cV\triangleq\Gamma(B_2)\diff(\Gamma(B_1)\cup\{B_1\})$, and
 $\cR\triangleq\cB\diff(\{B_1,B_2\}\cup\cM\cup\cU\cup\cV)$.}
 Write $\cF_{ij}\triangleq\{B\notin\{B_1,B_2\}:
 |B_1\cap B|=i,\ |B_2\cap B|=j\}$ for the fibres whose
 cardinalities are $x_{ij}$ in Lemma~\ref{hc:lem:common}.
 \step{hcLocalFibreM}{$\cM=\cF_{11}$.}
 \begin{proof}\pf\ Definition~\ref{hc:def:graph} identifies each
  adjacency with a one-point intersection and excludes the two
  endpoints by irreflexivity. Conversely the fibre gives both
  adjacencies. Lemma~\ref{hc:lem:intersections} places these
  intersections among the two possible sizes for distinct labels.~\qed\end{proof}
 \step{hcLocalFibreU}{$\cU=\cF_{12}$.}
 \begin{proof}\pf\ Membership in $\cU$ gives adjacency to $B_1$
  and excludes $B_2$ and adjacency to $B_2$.
  Definition~\ref{hc:def:graph} gives the first intersection size
  one; Lemma~\ref{hc:lem:intersections} forces the second to be two.
  Conversely these two sizes and the exclusions in $\cF_{12}$ give
  precisely the defining membership conditions of $\cU$.~\qed\end{proof}
 \step{hcLocalFibreV}{$\cV=\cF_{21}$.}
 \begin{proof}\pf\ Membership in $\cV$ gives adjacency to $B_2$
  and excludes $B_1$ and adjacency to $B_1$.
  Definition~\ref{hc:def:graph} gives the second size one;
  Lemma~\ref{hc:lem:intersections} forces the first to be two.
  Conversely the two sizes and endpoint exclusions in $\cF_{21}$
  give precisely the defining conditions of $\cV$.~\qed\end{proof}
 \step{hcLocalFibreR}{$\cR=\cF_{22}$.}
 \begin{proof}\pf\ Outside $\{B_1,B_2\}$,
  Lemma~\ref{hc:lem:intersections} leaves four possible pairs of
  sizes. Steps \stepref{hcLocalFibreM}--\stepref{hcLocalFibreV}
  remove the first three, leaving $(2,2)$. Conversely this pair
  gives no adjacency to either endpoint by Definition~\ref{hc:def:graph},
  so its label lies in none of $\cM,\cU,\cV$ and belongs to $\cR$.~\qed\end{proof}
 \step{hcLocalSizes}{$|\cM|=5$, $|\cU|=4$, $|\cV|=4$, and $|\cR|=6$.}
 \begin{proof}\pf\ Hypothesis \asm[2]{0} and
  Definition~\ref{hc:def:graph} give $s=1$.
  Apply the $s=1$ row of Lemma~\ref{hc:lem:common} to
  steps \stepref{hcLocalFibreM}--\stepref{hcLocalFibreR}.~\qed\end{proof}
 \step{hcLocalDisjoint}{$\cM,\cU,\cV,\cR$ are pairwise disjoint.
  Each is disjoint from $\{B_1,B_2\}$.}
 \begin{proof}\pf\ Their four different pairs of intersection
  sizes in steps \stepref{hcLocalFibreM}--\stepref{hcLocalFibreR}
  are mutually exclusive, and each fibre excludes the endpoints.~\qed\end{proof}
 \step{hcLocalCoverParts}{$\cM\cup\cU\cup\cV\cup\cR=
  \cB\diff\{B_1,B_2\}$.}
 \begin{proof}\pf\ By Lemma~\ref{hc:lem:intersections}, every label
  outside the endpoints has one of the four pairs of sizes.
  Use steps \stepref{hcLocalFibreM}--\stepref{hcLocalFibreR}.~\qed\end{proof}
 \step{hcLocalNeighboursOne}{$\Gamma(B_1)=\{B_2\}\cup\cM\cup\cU$,
  with pairwise disjoint parts.}
 \begin{proof}\pf\ The endpoint $B_2$ is a neighbour by \asm[2]{0}.
  Every other neighbour has first size one, so lies in $\cF_{11}$
  or $\cF_{12}$ by Lemma~\ref{hc:lem:intersections} and
  Definition~\ref{hc:def:graph}; both fibres consist of neighbours.
  Apply \stepref{hcLocalFibreM}, \stepref{hcLocalFibreU}, and
  \stepref{hcLocalDisjoint}.~\qed\end{proof}
 \step{hcLocalNeighboursTwo}{$\Gamma(B_2)=\{B_1\}\cup\cM\cup\cV$,
  with pairwise disjoint parts.}
 \begin{proof}\pf\ The endpoint $B_1$ is a neighbour by \asm[2]{0}.
  Every other neighbour has second size one, so lies in $\cF_{11}$
  or $\cF_{21}$ by Lemma~\ref{hc:lem:intersections} and
  Definition~\ref{hc:def:graph}; both fibres consist of neighbours.
  Apply \stepref{hcLocalFibreM}, \stepref{hcLocalFibreV}, and
  \stepref{hcLocalDisjoint}.~\qed\end{proof}
 \step{hcLocalU}{$\cU$ is a clique.}
 \begin{proof}
  \step{hcLocalUPair}{\assume{$Y,Y'\in\cU$, $Y\ne Y'$, and $Y\not\sim Y'$.}
   \prove{\textsc{False}.}}
  \begin{proof}\pf\ In \stepref{hcLocalFibreU}, $B_2,Y,Y'$ are
   distinct neighbours of $B_1$. Neither $Y$ nor $Y'$ is adjacent
   to $B_2$; \asm{2} excludes their mutual edge. This is the claw
   forbidden by Lemma~\ref{hc:lem:claw}.~\qed\end{proof}
  \qedstep
  \begin{proof}\pf\ Step \stepref{hcLocalUPair} excludes every
   missing edge between distinct members, as required by
   Definition~\ref{hc:def:graph}.~\qed\end{proof}
 \end{proof}
 \step{hcLocalV}{$\cV$ is a clique.}
 \begin{proof}
  \step{hcLocalVPair}{\assume{$W,W'\in\cV$, $W\ne W'$, and $W\not\sim W'$.}
   \prove{\textsc{False}.}}
  \begin{proof}\pf\ Step \stepref{hcLocalFibreV} makes
   $B_1,W,W'$ distinct neighbours of $B_2$; $B_1$ is adjacent to
   neither $W$ nor $W'$. With \asm{2} these are the three leaves
   of a claw, contradicting Lemma~\ref{hc:lem:claw}.~\qed\end{proof}
  \qedstep
  \begin{proof}\pf\ Step \stepref{hcLocalVPair} proves the clique
   predicate of Definition~\ref{hc:def:graph}.~\qed\end{proof}
 \end{proof}
 \step{hcLocalTwo}{For every $Y\in\cU$, $|\Gamma(Y)\cap\cM|=2$.}
 \begin{proof}
  \step{hcLocalYCommon}{For each $Y\in\cU$, the common neighbours of
   $B_1,Y$ are the disjoint union of $\cU\diff\{Y\}$ and
   $\Gamma(Y)\cap\cM$.}
  \begin{proof}\pf\ In the first neighbourhood decomposition of
   \stepref{hcLocalNeighboursOne}, $B_2\not\sim Y$ by membership in
   $\cU$; step \stepref{hcLocalU} gives every member of
   $\cU\diff\{Y\}$. These leave precisely the stated part of $\cM$.~\qed\end{proof}
  \qedstep
  \begin{proof}\pf\ Lemma~\ref{hc:lem:common} gives five common
   neighbours of the adjacent pair $B_1,Y$. The first part in
   \stepref{hcLocalYCommon} has size $4-1=3$ by
   \stepref{hcLocalSizes} and \ref{hc:facts:finite}, leaving two.~\qed\end{proof}
 \end{proof}
 \step{hcLocalEdgeSwap}{$\displaystyle
  \sum_{X\in\cM}|\Gamma(X)\cap\cU|
  =\sum_{Y\in\cU}|\Gamma(Y)\cap\cM|$.}
 \begin{proof}\pf\ Write each cardinality as a sum of adjacency
  indicators. Interchange the finite sums by \path{Finset.sum_comm}
  in \ref{hc:facts:finite}; symmetry of adjacency from
  Definition~\ref{hc:def:graph} gives the other sum.~\qed\end{proof}
 \step{hcLocalEight}{$\displaystyle
  \sum_{X\in\cM}|\Gamma(X)\cap\cU|=8$.}
 \begin{proof}\pf\ By \stepref{hcLocalEdgeSwap} the sum equals
  $\sum_{Y\in\cU}2$, using \stepref{hcLocalTwo}.
  Step \stepref{hcLocalSizes} gives $4\cdot2=8$.~\qed\end{proof}
 \step{hcLocalPick}{\pick\ $B_3\in\cM$ such that $|\Gamma(B_3)\cap\cU|\ge2$.}
 \begin{proof}\pf\ Otherwise each summand of \stepref{hcLocalEight}
  is at most one. Termwise summation gives
  $8\le\sum_{X\in\cM}1=|\cM|=5$, using
  \stepref{hcLocalSizes} and \ref{hc:facts:finite}, a contradiction.~\qed\end{proof}
 \step{hcLocalDefect}{\pflet{$j\triangleq|\cU\diff\Gamma(B_3)|$, and}
  \[
  h\triangleq|\Gamma(B_3)\cap(\cM\diff\{B_3\})|,\quad
  e\triangleq|\Gamma(B_3)\cap\cV|,\quad
  \ell\triangleq|\Gamma(B_3)\cap\cR|.
  \]
  Then $0\le j\le2$ and $|\Gamma(B_3)\cap\cU|=4-j$.}
 \begin{proof}\pf\ The neighbours and nonneighbours of $B_3$ partition
  the four-element set $\cU$, by \stepref{hcLocalSizes} and
  \ref{hc:facts:finite}; \stepref{hcLocalPick} supplies at least two neighbours.~\qed\end{proof}
 \step{hcLocalCounts}{$h=j$, $e=4-j$, and $\ell=j$.}
 \begin{proof}
  \step{hcLocalFirst}{$1+h+(4-j)=5$.}
  \begin{proof}\pf\ Count common neighbours of $B_1,B_3$ using
   Lemma~\ref{hc:lem:common} and \stepref{hcLocalNeighboursOne}:
   $B_2$ contributes one, $\cM\diff\{B_3\}$ contributes $h$, and
   $\cU$ contributes $4-j$ by \stepref{hcLocalDefect}.~\qed\end{proof}
  \step{hcLocalSecond}{$1+h+e=5$.}
  \begin{proof}\pf\ Count common neighbours of $B_2,B_3$ using
   Lemma~\ref{hc:lem:common}: $B_1$ contributes one, the part
   $\cM\diff\{B_3\}$ contributes $h$, and $\cV$ contributes $e$,
   by \stepref{hcLocalNeighboursTwo} and \stepref{hcLocalDefect}.~\qed\end{proof}
  \step{hcLocalDegree}{$2+h+(4-j)+e+\ell=10$.}
  \begin{proof}\pf\ Count all neighbours of $B_3$ in the partition
   \stepref{hcLocalCoverParts}, whose parts are disjoint by
   \stepref{hcLocalDisjoint}. Both $B_1,B_2$ are neighbours by
   \stepref{hcLocalPick}; the other contributions are the four
   counts in \stepref{hcLocalDefect}. The degree is ten by
   Lemma~\ref{hc:lem:intersections}.~\qed\end{proof}
  \qedstep
  \begin{proof}\pf\ Solve \stepref{hcLocalFirst} for $h$, then
   \stepref{hcLocalSecond} for $e$, then \stepref{hcLocalDegree}
   for $\ell$, over $\ints$.~\qed\end{proof}
 \end{proof}
 \step{hcLocalParts}{\pflet{$\cU_0\triangleq\cU\diff\Gamma(B_3)$ and
  $\cV_1\triangleq\cV\cap\Gamma(B_3)$.}
  They are disjoint cliques of sizes $j,4-j$, disjoint also from
  $\{B_1,B_2,B_3\}$.}
 \begin{proof}\pf\ Disjointness is \stepref{hcLocalDisjoint} and
  \stepref{hcLocalPick}; the sizes are \stepref{hcLocalDefect} and
  \stepref{hcLocalCounts}. They are subsets of the cliques in
  \stepref{hcLocalU} and \stepref{hcLocalV}.~\qed\end{proof}
 \step{hcLocalNoCross}{No member of $\cU_0$ is adjacent to a member of $\cV_1$.}
 \begin{proof}
  \step{hcLocalCrossData}{\assume{$Y\in\cU_0$ and $W\in\cV_1$.}
   \prove{$Y\not\sim W$.}}
  \begin{proof}
   \step{hcLocalFiveDistinct}{$B_1,B_2,B_3,Y,W$ are distinct.}
   \begin{proof}\pf\ Step \stepref{hcLocalParts} gives disjointness,
    and \stepref{hcLocalDisjoint}, \stepref{hcLocalPick} separate
    $B_3$ from the adjacent pair $B_1,B_2$.~\qed\end{proof}
   \step{hcLocalFiveTriangle}{$B_1,B_2,B_3$ form a triangle.}
   \begin{proof}\pf\ Use \asm[2]{0} and $B_3\in\cM$ in
    \stepref{hcLocalPick}.~\qed\end{proof}
   \step{hcLocalFiveY}{$Y\sim B_1$, $Y\not\sim B_2$, and $Y\not\sim B_3$.}
   \begin{proof}\pf\ Membership in $\cU_0$ in \asm{2}, and the
    definition in \stepref{hcLocalParts} and the fibre description
    in \stepref{hcLocalFibreU}.~\qed\end{proof}
   \step{hcLocalFiveW}{$W\sim B_2$, $W\sim B_3$, and $W\not\sim B_1$.}
   \begin{proof}\pf\ Membership in $\cV_1$ in \asm{2}, and the
    definition in \stepref{hcLocalParts} and the fibre description
    in \stepref{hcLocalFibreV}.~\qed\end{proof}
   \qedstep
   \begin{proof}\pf\ Apply Lemma~\ref{hc:lem:five} to
    \stepref{hcLocalFiveDistinct}--\stepref{hcLocalFiveW}.~\qed\end{proof}
  \end{proof}
  \qedstep
  \begin{proof}\pf\ Step \stepref{hcLocalCrossData} is quantified
   over every pair of the indicated parts.~\qed\end{proof}
 \end{proof}
 \step{hcLocalPattern}{On the named blocks $B_1,B_2,B_3$ and the
  members of $\cU_0,\cV_1$, $C$ is symmetric with diagonal $10$.
  Its off-diagonal entries are
  \[
  \begin{array}{c|ccc|cc}
       &B_1&B_2&B_3&\cU_0&\cV_1\\\hline
  B_1&*&3&3&3&-4\\
  B_2&3&*&3&-4&3\\
  B_3&3&3&*&-4&3\\\hline
  \cU_0&3&-4&-4&3&-4\\
  \cV_1&-4&3&3&-4&3
  \end{array}
  \]
  A set in this table denotes any of its members, with distinct
  blocks required for within-set entries.}
 \begin{proof}
  \step{hcPatternDiagonal}{For every block $X$, $C(X,X)=10$.}
  \begin{proof}\pf\ Lemma~\ref{hc:lem:Connor}.~\qed\end{proof}
  \step{hcPatternTriangle}{The three entries among $B_1,B_2,B_3$ are $3$.}
  \begin{proof}\pf\ The triangle follows from \asm[2]{0},
   \stepref{hcLocalFibreM}, and \stepref{hcLocalPick};
   use Lemma~\ref{hc:lem:Connor}.~\qed\end{proof}
  \step{hcPatternOneU}{For $Y\in\cU_0$, $C(B_1,Y)=3$.}
  \begin{proof}\pf\ Membership in $\cU$ gives adjacency to $B_1$
   by \stepref{hcLocalFibreU} and \stepref{hcLocalParts};
   apply Lemma~\ref{hc:lem:Connor}.~\qed\end{proof}
  \step{hcPatternTwoU}{For $Y\in\cU_0$, $C(B_2,Y)=-4$.}
  \begin{proof}\pf\ Step \stepref{hcLocalFibreU} excludes adjacency
   to $B_2$; distinctness in \stepref{hcLocalParts} permits
   Lemma~\ref{hc:lem:Connor}.~\qed\end{proof}
  \step{hcPatternThreeU}{For $Y\in\cU_0$, $C(B_3,Y)=-4$.}
  \begin{proof}\pf\ The definition and disjointness in
   \stepref{hcLocalParts} give nonadjacency and distinctness;
   apply Lemma~\ref{hc:lem:Connor}.~\qed\end{proof}
  \step{hcPatternOneV}{For $W\in\cV_1$, $C(B_1,W)=-4$.}
  \begin{proof}\pf\ Step \stepref{hcLocalFibreV} excludes adjacency
   to $B_1$; use the distinctness in \stepref{hcLocalParts} and
   Lemma~\ref{hc:lem:Connor}.~\qed\end{proof}
  \step{hcPatternTwoV}{For $W\in\cV_1$, $C(B_2,W)=3$.}
  \begin{proof}\pf\ Step \stepref{hcLocalFibreV} gives adjacency to
   $B_2$ for every member of the subset in \stepref{hcLocalParts};
   use Lemma~\ref{hc:lem:Connor}.~\qed\end{proof}
  \step{hcPatternThreeV}{For $W\in\cV_1$, $C(B_3,W)=3$.}
  \begin{proof}\pf\ The definition of $\cV_1$ in
   \stepref{hcLocalParts} gives adjacency; use Lemma~\ref{hc:lem:Connor}.~\qed\end{proof}
  \step{hcPatternUU}{For distinct $Y,Y'\in\cU_0$, $C(Y,Y')=3$.}
  \begin{proof}\pf\ The clique in \stepref{hcLocalParts} gives
   adjacency, so Lemma~\ref{hc:lem:Connor} applies.~\qed\end{proof}
  \step{hcPatternVV}{For distinct $W,W'\in\cV_1$, $C(W,W')=3$.}
  \begin{proof}\pf\ The clique in \stepref{hcLocalParts} gives
   adjacency; use Lemma~\ref{hc:lem:Connor}.~\qed\end{proof}
  \step{hcPatternUV}{For $Y\in\cU_0$ and $W\in\cV_1$, $C(Y,W)=-4$.}
  \begin{proof}\pf\ Step \stepref{hcLocalNoCross} excludes the edge;
   \stepref{hcLocalParts} gives distinctness, so
   Lemma~\ref{hc:lem:Connor} applies.~\qed\end{proof}
  \qedstep
  \begin{proof}\pf\ Steps \stepref{hcPatternDiagonal}--\stepref{hcPatternUV}
   cover the diagonal and all ten unordered pair types.
   Symmetry from Lemma~\ref{hc:lem:Connor} gives the reversed entries.~\qed\end{proof}
 \end{proof}
 \step{hcLocalNotOne}{\assume{$j=1$.}\prove{\textsc{False}.}}
 \begin{proof}
  \step{hcLocalEnumerateOne}{\pick\ $Y,W_1,W_2,W_3$ with
   $\cU_0=\{Y\}$, $\cV_1=\{W_1,W_2,W_3\}$, and the $W_i$ pairwise distinct.}
  \begin{proof}\pf\ By \asm{1} and \stepref{hcLocalParts}, the sizes
   are one and three. Apply \path{Finset.card_eq_one} and
   \path{Finset.card_eq_three} from \ref{hc:facts:finite}.~\qed\end{proof}
  \step{hcLocalInjectionOne}{Put
   $a\triangleq(B_1,B_2,B_3,Y,W_1,W_2,W_3):\mathrm{Fin}\,7\to\cB$.
   Then $a$ is injective.}
  \begin{proof}\pf\ The first three blocks are distinct by
   \asm[2]{0}, \stepref{hcLocalDisjoint}, and \stepref{hcLocalPick}.
   Step \stepref{hcLocalParts} separates them from the two disjoint
   parts, and \stepref{hcLocalEnumerateOne} separates the $W_i$.
   Thus the literal Lean vector \texttt{![B1, B2, B3, Y, W1, W2, W3]}
   has seven distinct entries.~\qed\end{proof}
  \step{hcLocalQOne}{$C_a=Q_1$.}
  \begin{proof}\pf\ Use the fixed list in \stepref{hcLocalInjectionOne}.
   In zero-based indices, the named blocks occupy $0,1,2$, $Y$ occupies
   $3$, and $W_1,W_2,W_3$ occupy $4,5,6$ by \stepref{hcLocalEnumerateOne}.
   For each of the $49$ pairs, the corresponding entry of
   \stepref{hcLocalPattern} equals the displayed entry of $Q_1$ in
   Definition~\ref{hc:def:certificates}: the diagonal is $10$,
   triangle and within-$W$ entries are $3$, the three entries from
   $Y$ to $B_1,B_2,B_3$ are $(3,-4,-4)$, each $W_i$ has
   $(-4,3,3)$ to those blocks, and all $Y,W_i$ entries are $-4$.
   These cases and their reverses exhaust the entries.~\qed\end{proof}
  \step{hcLocalSingularOne}{$\det C_a=0$.}
  \begin{proof}\pf\ Apply Lemma~\ref{hc:lem:principal} to
   \stepref{hcLocalInjectionOne}, whose domain has size seven.~\qed\end{proof}
  \qedstep
  \begin{proof}\pf\ Steps \stepref{hcLocalQOne} and
   \stepref{hcLocalSingularOne} contradict Lemma~\ref{hc:lem:inverse} at $j=1$.~\qed\end{proof}
 \end{proof}
 \step{hcLocalNotTwo}{\assume{$j=2$.}\prove{\textsc{False}.}}
 \begin{proof}
  \step{hcLocalEnumerateTwo}{\pick\ $Y_1,Y_2,W_1,W_2$ with
   $\cU_0=\{Y_1,Y_2\}$, $\cV_1=\{W_1,W_2\}$,
   $Y_1\ne Y_2$, and $W_1\ne W_2$.}
  \begin{proof}\pf\ The two sizes in \stepref{hcLocalParts} are
   both two by \asm{1}. Apply \path{Finset.card_eq_two} from
   \ref{hc:facts:finite} to each set.~\qed\end{proof}
  \step{hcLocalInjectionTwo}{Put
   $a\triangleq(B_1,B_2,B_3,Y_1,Y_2,W_1,W_2):\mathrm{Fin}\,7\to\cB$.
   Then $a$ is injective.}
  \begin{proof}\pf\ The initial blocks are distinct by
   \asm[2]{0}, \stepref{hcLocalDisjoint}, and \stepref{hcLocalPick}.
   The disjointness in \stepref{hcLocalParts} and the two inequalities
   in \stepref{hcLocalEnumerateTwo} separate all remaining pairs.
   This is the literal Lean vector \texttt{![B1, B2, B3, Y1, Y2, W1, W2]}.~\qed\end{proof}
  \step{hcLocalQTwo}{$C_a=Q_2$.}
  \begin{proof}\pf\ The list in \stepref{hcLocalInjectionTwo} puts
   $B_1,B_2,B_3$ at indices $0,1,2$, $Y_1,Y_2$ at $3,4$, and
   $W_1,W_2$ at $5,6$, using \stepref{hcLocalEnumerateTwo}.
   For each of the $49$ pairs, \stepref{hcLocalPattern} matches
   $Q_2$ in Definition~\ref{hc:def:certificates}: the diagonal is
   $10$, the triangle, $Y_1,Y_2$, and $W_1,W_2$ entries are $3$,
   each $Y_i$ has $(3,-4,-4)$ to the triangle, each $W_i$ has
   $(-4,3,3)$, and the four $Y_i,W_h$ entries are $-4$.
   The reverse entries agree by symmetry.~\qed\end{proof}
  \step{hcLocalSingularTwo}{$\det C_a=0$.}
  \begin{proof}\pf\ Lemma~\ref{hc:lem:principal} applies to
   \stepref{hcLocalInjectionTwo} with domain size seven.~\qed\end{proof}
  \qedstep
  \begin{proof}\pf\ Steps \stepref{hcLocalQTwo} and
   \stepref{hcLocalSingularTwo} contradict Lemma~\ref{hc:lem:inverse} at $j=2$.~\qed\end{proof}
 \end{proof}
 \step{hcLocalZero}{$j=0$; $B_3$ is adjacent to every member of
  $\cU$ and to no member of $\cM\diff\{B_3\}$.}
 \begin{proof}\pf\ The integer bound $0\le j\le2$ of
  \stepref{hcLocalDefect} and the exclusions
  \stepref{hcLocalNotOne}, \stepref{hcLocalNotTwo} force zero.
  The counts in \stepref{hcLocalDefect}, \stepref{hcLocalCounts}
  are then four and zero; use \ref{hc:facts:finite}.~\qed\end{proof}
 \step{hcLocalM}{$\cM\diff\{B_3\}$ is a clique of size four.}
 \begin{proof}
  \step{hcLocalMSize}{$|\cM\diff\{B_3\}|=4$.}
  \begin{proof}\pf\ Use $|\cM|=5$, $B_3\in\cM$ from
   \stepref{hcLocalSizes}, \stepref{hcLocalPick}, and
   deletion of a singleton in \ref{hc:facts:finite}.~\qed\end{proof}
  \step{hcLocalMClaw}{\assume{$X,Y\in\cM\diff\{B_3\}$ are distinct and nonadjacent.}
   \prove{\textsc{False}.}}
  \begin{proof}\pf\ The three distinct vertices $B_3,X,Y$ are
   neighbours of $B_1$ by \stepref{hcLocalFibreM} and
   \stepref{hcLocalPick}; $B_3$ has no edges to $X,Y$ by
   \stepref{hcLocalZero}. With \asm{2} they give a claw,
   contradicting Lemma~\ref{hc:lem:claw}.~\qed\end{proof}
  \qedstep
  \begin{proof}\pf\ Step \stepref{hcLocalMSize} gives the size and
   \stepref{hcLocalMClaw} gives every required edge.~\qed\end{proof}
 \end{proof}
 \step{hcLocalSOne}{Put $S_1\triangleq\{B_1,B_2\}\cup(\cM\diff\{B_3\})$.
  Then $S_1\in\cS$ and $B_1,B_2\in S_1$.}
 \begin{proof}
  \step{hcLocalSOneSize}{$|S_1|=6$.}
  \begin{proof}\pf\ The union is disjoint by
   \stepref{hcLocalDisjoint}, and \stepref{hcLocalM} gives $2+4=6$.~\qed\end{proof}
  \step{hcLocalSOneCore}{$B_1,B_2$ are adjacent.}
  \begin{proof}\pf\ Hypothesis \asm[2]{0}.~\qed\end{proof}
  \step{hcLocalSOneMixed}{Every member of $\cM\diff\{B_3\}$ is
   adjacent to each of $B_1,B_2$.}
  \begin{proof}\pf\ The definition of $\cM$ in \stepref{hcLocalFibreM}.~\qed\end{proof}
  \step{hcLocalSOneInternal}{Distinct members of $\cM\diff\{B_3\}$ are adjacent.}
  \begin{proof}\pf\ Step \stepref{hcLocalM}.~\qed\end{proof}
  \qedstep
  \begin{proof}\pf\ The three pair types in
   \stepref{hcLocalSOneCore}--\stepref{hcLocalSOneInternal} exhaust
   pairs in $S_1$. With \stepref{hcLocalSOneSize} this is membership
   in $\cS$, by Definition~\ref{hc:def:graph}; the two containments
   follow from the displayed union.~\qed\end{proof}
 \end{proof}
 \step{hcLocalSTwo}{Put $S_2\triangleq\{B_1,B_3\}\cup\cU$.
  Then $S_2\in\cS$.}
 \begin{proof}
  \step{hcLocalSTwoSize}{$|S_2|=6$.}
  \begin{proof}\pf\ The two-element set is disjoint from $\cU$
   by \stepref{hcLocalDisjoint} and \stepref{hcLocalPick}; its
   union has size $2+4=6$ by \stepref{hcLocalSizes}.~\qed\end{proof}
  \step{hcLocalSTwoCore}{$B_1\sim B_3$.}
  \begin{proof}\pf\ Step \stepref{hcLocalPick} puts $B_3\in\cM$.~\qed\end{proof}
  \step{hcLocalSTwoFirst}{Every member of $\cU$ is adjacent to $B_1$.}
  \begin{proof}\pf\ Step \stepref{hcLocalFibreU}.~\qed\end{proof}
  \step{hcLocalSTwoThird}{Every member of $\cU$ is adjacent to $B_3$.}
  \begin{proof}\pf\ Step \stepref{hcLocalZero}.~\qed\end{proof}
  \step{hcLocalSTwoInternal}{Every pair of distinct members of $\cU$ is adjacent.}
  \begin{proof}\pf\ Step \stepref{hcLocalU}.~\qed\end{proof}
  \qedstep
  \begin{proof}\pf\ Steps \stepref{hcLocalSTwoCore}--\stepref{hcLocalSTwoInternal}
   cover all pair types. Together with \stepref{hcLocalSTwoSize},
   they are the predicate of Definition~\ref{hc:def:graph}.~\qed\end{proof}
 \end{proof}
 \step{hcLocalMeet}{$S_1\cap S_2=\{B_1\}$.}
 \begin{proof}\pf\ In \stepref{hcLocalSOne} and \stepref{hcLocalSTwo},
  the parts $\{B_2\}$, $\cM\diff\{B_3\}$, $\{B_3\}$, and
  $\cU$ are pairwise disjoint by \stepref{hcLocalDisjoint} and
  \stepref{hcLocalPick}; only $B_1$ is shared.~\qed\end{proof}
 \step{hcLocalCover}{$S_1\cup S_2=\{B_1\}\cup\Gamma(B_1)$.}
 \begin{proof}\pf\ The two unions in \stepref{hcLocalSOne},
  \stepref{hcLocalSTwo} combine to $\{B_1,B_2\}\cup\cM\cup\cU$,
  since $B_3\in\cM$. Use the first neighbourhood decomposition
  of \stepref{hcLocalNeighboursOne}.~\qed\end{proof}
 \step{hcLocalUnique}{\assume{$K\in\cS$ and $B_1,B_2\in K$.}
  \prove{$K=S_1$.}}
 \begin{proof}
  \step{hcLocalKFour}{$K\diff\{B_1,B_2\}\subseteq\cM$ and
   $|K\diff\{B_1,B_2\}|=4$.}
  \begin{proof}\pf\ Every remaining vertex is adjacent to both
   initial blocks by Definition~\ref{hc:def:graph}. The two
   distinct blocks are in the six-element set $K$, so
   \ref{hc:facts:finite} gives size four.~\qed\end{proof}
  \step{hcLocalKExclude}{$B_3\notin K$.}
  \begin{proof}
   \step{hcLocalKBad}{\assume{$B_3\in K$.}\prove{\textsc{False}.}}
   \begin{proof}\pf\ By \stepref{hcLocalKFour}, removing $B_3$
    leaves three vertices in $\cM\diff\{B_3\}$, thus at least
    one, by \ref{hc:facts:finite}. The clique $K$ makes it adjacent
    to $B_3$, contradicting \stepref{hcLocalZero}.~\qed\end{proof}
   \qedstep
   \begin{proof}\pf\ Step \stepref{hcLocalKBad} proves the negation.~\qed\end{proof}
  \end{proof}
  \step{hcLocalKEqual}{$K\diff\{B_1,B_2\}=\cM\diff\{B_3\}$.}
  \begin{proof}\pf\ Steps \stepref{hcLocalKFour}, \stepref{hcLocalKExclude}
   give inclusion; both sets have size four by \stepref{hcLocalM}.
   Apply the subset and cardinality criterion in \ref{hc:facts:finite}.~\qed\end{proof}
  \qedstep
  \begin{proof}\pf\ Reinsert the two members $B_1,B_2$ into
   \stepref{hcLocalKEqual}; the resulting union is $S_1$ by
   \stepref{hcLocalSOne}.~\qed\end{proof}
 \end{proof}
 \qedstep
 \begin{proof}\pf\ Steps \stepref{hcLocalSOne}, \stepref{hcLocalSTwo},
  \stepref{hcLocalMeet}, \stepref{hcLocalCover}, and
  \stepref{hcLocalUnique} give the stated witnesses and uniqueness.~\qed\end{proof}
\end{proof}

\begin{namedresult}{Corollary}{(the unique six-clique through an edge)}\label{hc:cor:unique}
 \lean{Design.exists_unique_sixClique_of_adj}\lean{Design.sixClique_eq_of_adj}
 \uses{hc:def:design,hc:def:graph,hc:lem:split}
 \leavevmode\par
 \assume{\begin{pfenum}
  \item $P,\cB$ are finite types with decidable equality and
  $D$ is a design with parameters $(15,5,21,7,2)$.
  \item $B_1\sim B_2$.
 \end{pfenum}}
 \prove{There is a unique $S\in\cS$ containing $B_1,B_2$.
 In particular, if $K,L\in\cS$ both contain $B_1,B_2$, then $K=L$.}
\end{namedresult}
\begin{proof}
 \pfsketch\ Choose one local witness and compare every other
 six-clique through the edge with that same witness.
 \step{hcUniqueWitness}{\pick\ $S_1,S_2$ with the conclusions of
  Lemma~\ref{hc:lem:split} for $B_1,B_2$.}
 \begin{proof}\pf\ Apply that lemma to \asm[2]{0}.~\qed\end{proof}
 \step{hcUniqueExists}{There is a unique $S\in\cS$ containing $B_1,B_2$.}
 \begin{proof}\pf\ Step \stepref{hcUniqueWitness} gives the witness
  $S_1$ and states that every such $S$ equals $S_1$.~\qed\end{proof}
 \step{hcUniqueEqual}{\assume{$K,L\in\cS$ and $B_1,B_2\in K\cap L$.}
  \prove{$K=L$.}}
 \begin{proof}\pf\ The uniqueness clause of \stepref{hcUniqueWitness}
  gives $K=S_1$ and $L=S_1$. Transitivity gives $K=L$.~\qed\end{proof}
 \qedstep
 \begin{proof}\pf\ Steps \stepref{hcUniqueExists} and \stepref{hcUniqueEqual}.~\qed\end{proof}
\end{proof}

\begin{namedresult}{Corollary}{(exactly two six-cliques through a block)}\label{hc:cor:two}
 \lean{Design.card_sixCliques_through_block}
 \uses{hc:def:design,hc:def:graph,hc:lem:intersections,hc:lem:split,hc:cor:unique,hc:facts:finite}
 \leavevmode\par
 \assume{\begin{pfenum}
  \item $P,\cB$ are finite types with decidable equality and
  $D$ is a design with parameters $(15,5,21,7,2)$.
  \item $B_1\in\cB$.
 \end{pfenum}}
 \prove{There are distinct $S_1,S_2\in\cS$ with
 $\{S\in\cS:B_1\in S\}=\{S_1,S_2\}$ and
 $S_1\cap S_2=\{B_1\}$. In particular the filtered family has size two.}
\end{namedresult}
\begin{proof}
 \pfsketch\ We will use one local split, then apply the stated
 adjacent-pair uniqueness lemma to each possible six-clique.
 \step{hcTwoPick}{\pick\ $B_2\in\Gamma(B_1)$.}
 \begin{proof}\pf\ Lemma~\ref{hc:lem:intersections} gives ten
  neighbours; positive cardinality in \ref{hc:facts:finite} gives a witness.~\qed\end{proof}
 \step{hcTwoSplit}{\pick\ $S_1,S_2$ with all the conclusions of
  Lemma~\ref{hc:lem:split} for $B_1,B_2$.}
 \begin{proof}\pf\ Apply that lemma using \stepref{hcTwoPick}.~\qed\end{proof}
 \step{hcTwoDistinct}{$S_1\ne S_2$, and both contain $B_1$.}
 \begin{proof}\pf\ Their intersection is $\{B_1\}$ by
  \stepref{hcTwoSplit}. Equality would give a singleton of size
  six, contrary to Definition~\ref{hc:def:graph}.~\qed\end{proof}
 \step{hcTwoExhaust}{\assume{$K\in\cS$ and $B_1\in K$.}
  \prove{$K=S_1$ or $K=S_2$.}}
 \begin{proof}
  \step{hcTwoOther}{\pick\ $X\in K\diff\{B_1\}$.}
  \begin{proof}\pf\ This set has size five by \asm{1},
   Definition~\ref{hc:def:graph}, and \ref{hc:facts:finite}.~\qed\end{proof}
  \step{hcTwoEdge}{$B_1\sim X$ and
   $X\in(S_1\diff\{B_1\})\cup(S_2\diff\{B_1\})$.}
  \begin{proof}\pf\ The clique $K$ in \asm{1} and
   \stepref{hcTwoOther} give adjacency. The covering conclusion
   of \stepref{hcTwoSplit} gives the indicated membership.~\qed\end{proof}
  \step{hcTwoFirst}{\case{$X\in S_1$.} $K=S_1$.}
  \begin{proof}\pf\ Both six-cliques contain the adjacent pair
   $B_1,X$ by \asm{1}, \stepref{hcTwoDistinct}, and
   \stepref{hcTwoEdge}. Apply the uniqueness conclusion of
   Corollary~\ref{hc:cor:unique} to this pair.~\qed\end{proof}
  \step{hcTwoSecond}{\case{$X\in S_2$.} $K=S_2$.}
  \begin{proof}\pf\ The cliques $K,S_2$ contain $B_1,X$ by
   \asm{1}, \stepref{hcTwoDistinct}, and \stepref{hcTwoEdge}.
   Corollary~\ref{hc:cor:unique} applied to $B_1,X$ identifies them.~\qed\end{proof}
  \qedstep
  \begin{proof}\pf\ Step \stepref{hcTwoEdge} makes
   \stepref{hcTwoFirst}, \stepref{hcTwoSecond} exhaustive.~\qed\end{proof}
 \end{proof}
 \qedstep
 \begin{proof}\pf\ Step \stepref{hcTwoExhaust} gives containment
  of the filtered family in $\{S_1,S_2\}$; \stepref{hcTwoDistinct}
  gives the reverse containment and size two. The intersection
  is part of \stepref{hcTwoSplit}.~\qed\end{proof}
\end{proof}

\subsubsection{The global family}

\begin{namedresult}{Lemma}{(seven six-cliques)}\label{hc:lem:seven}
 \lean{Design.card_sixCliques}
 \uses{hc:def:design,hc:def:graph,hc:cor:two,hc:facts:finite}
 \leavevmode\par
 \assume{\begin{pfenum}
  \item $P,\cB$ are finite types with decidable equality.
  \item $D$ is a design with parameters $(15,5,21,7,2)$.
 \end{pfenum}}
 \prove{$|\cS|=7$.}
\end{namedresult}
\begin{proof}
 \pfsketch\ Count incidences between block labels and six-cliques.
 \step{hcSevenByClique}{There are $6|\cS|$ pairs $(B,S)$ with $B\in S\in\cS$.}
 \begin{proof}\pf\ Every member of $\cS$ has size six by
  Definition~\ref{hc:def:graph}; apply \ref{hc:facts:finite}.~\qed\end{proof}
 \step{hcSevenByBlock}{There are $42$ such pairs.}
 \begin{proof}\pf\ Each of the $21$ block labels belongs to two
  members of $\cS$ by Corollary~\ref{hc:cor:two} and
  Definition~\ref{hc:def:design}; use \ref{hc:facts:finite}.~\qed\end{proof}
 \qedstep
 \begin{proof}\pf\ Equate \stepref{hcSevenByClique} and
  \stepref{hcSevenByBlock}, then divide $6|\cS|=42$ by six.~\qed\end{proof}
\end{proof}

\begin{namedresult}{Lemma}{(intersections of six-cliques)}\label{hc:lem:clique-meet}
 \lean{Design.sixClique_inter_card}
 \uses{hc:def:design,hc:def:graph,hc:cor:unique,hc:cor:two,hc:lem:seven,hc:facts:finite}
 \leavevmode\par
 \assume{\begin{pfenum}
  \item $P,\cB$ are finite types with decidable equality and
  $D$ is a design with parameters $(15,5,21,7,2)$.
  \item $S_0,S_1\in\cS$ and $S_0\ne S_1$.
 \end{pfenum}}
 \prove{$|S_0\cap S_1|=1$.}
\end{namedresult}
\begin{proof}
 \pfsketch\ Any two distinct cliques share at most one block.
 The other clique through each of the six blocks of $S_0$ then
 enumerates all six cliques different from $S_0$.
 \step{hcMeetAtMost}{For all distinct $S,T\in\cS$, $|S\cap T|\le1$.}
 \begin{proof}
  \step{hcMeetTwo}{\assume{$S,T\in\cS$, $S\ne T$, and distinct $B,B'\in S\cap T$.}
   \prove{\textsc{False}.}}
  \begin{proof}\pf\ $B\sim B'$ by the clique predicate in
   Definition~\ref{hc:def:graph}. Corollary~\ref{hc:cor:unique}
   identifies $S,T$ through this edge, contradicting $S\ne T$.~\qed\end{proof}
  \qedstep
  \begin{proof}\pf\ The cardinality-at-most-one criterion of
   \ref{hc:facts:finite} applies to \stepref{hcMeetTwo}.~\qed\end{proof}
 \end{proof}
 \step{hcMeetMap}{For every $B\in S_0$ there is a unique
  $f(B)\in\cS\diff\{S_0\}$ containing $B$.}
 \begin{proof}\pf\ Corollary~\ref{hc:cor:two} supplies exactly
  two cliques through $B$, one of which is $S_0$ by \asm[2]{0}.
  Deleting that one leaves a singleton by \ref{hc:facts:finite}.~\qed\end{proof}
 \step{hcMeetInjective}{$f:S_0\to\cS\diff\{S_0\}$ is injective.}
 \begin{proof}\pf\ If $f(B)=f(B')$ with $B\ne B'$, then
  $S_0$ and $f(B)\ne S_0$ share two blocks by \stepref{hcMeetMap},
  contradicting \stepref{hcMeetAtMost}.~\qed\end{proof}
 \step{hcMeetSurjective}{$f$ is surjective.}
 \begin{proof}\pf\ Its domain has size six by
  Definition~\ref{hc:def:graph}; its codomain has size $7-1=6$
  by Lemma~\ref{hc:lem:seven} and \asm[2]{0}.
  The finite bijectivity criterion of \ref{hc:facts:finite}
  applies to \stepref{hcMeetInjective}.~\qed\end{proof}
 \step{hcMeetNonempty}{$S_0\cap S_1$ is nonempty.}
 \begin{proof}\pf\ By \asm[2]{0} and \stepref{hcMeetSurjective},
  choose $B\in S_0$ with $f(B)=S_1$. Step \stepref{hcMeetMap}
  gives $B\in S_1$.~\qed\end{proof}
 \qedstep
 \begin{proof}\pf\ Combine the lower bound from
  \stepref{hcMeetNonempty} with \stepref{hcMeetAtMost}, using
  \ref{hc:facts:finite}.~\qed\end{proof}
\end{proof}

\begin{namedresult}{Lemma}{(points in a six-clique)}\label{hc:lem:point-clique}
 \lean{Design.card_blocksThrough_in_sixClique}
 \uses{hc:def:design,hc:def:graph,hc:facts:finite}
 \leavevmode\par
 \assume{\begin{pfenum}
  \item $P,\cB$ are finite types with decidable equality and
  $D$ is a design with parameters $(15,5,21,7,2)$.
  \item $S_0\in\cS$.
 \end{pfenum}}
 \prove{For every $p\in P$, $|\{B\in S_0:p\in B\}|=2$.}
\end{namedresult}
\begin{proof}
 \pfsketch\ Two moments of the point multiplicities give a sum
 of squares equal to zero.
 \step{hcPointDefine}{\pflet{$m(p)\triangleq|\{B\in S_0:p\in B\}|$.}
  Then $\sum_{p\in P}m(p)=30$.}
 \begin{proof}\pf\ Count $(p,B)$ with $B\in S_0$ and $p\in B$.
  There are six labels, each with five points, by
  Definitions~\ref{hc:def:graph}, \ref{hc:def:design}, and
  \ref{hc:facts:finite}.~\qed\end{proof}
 \step{hcPointSecond}{$\sum_{p\in P}m(p)(m(p)-1)=30$ in $\ints$.}
 \begin{proof}\pf\ Count $(p,B,B')$ with distinct $B,B'\in S_0$
  and $p\in B\cap B'$. Each of the $6\cdot5$ ordered pairs
  meets in exactly one point by Definition~\ref{hc:def:graph}.
  Counting by $p$ gives the sum, using \ref{hc:facts:finite}.~\qed\end{proof}
 \step{hcPointSquares}{$\sum_{p\in P}(m(p)-2)^2=0$ in $\ints$.}
 \begin{proof}
  \pf\ Integer expansion gives
  \begin{align*}
   \sum_p(m(p)-2)^2
    &=\sum_p m(p)(m(p)-1)-3\sum_p m(p)+4|P|
        &&\text{[ring arithmetic]}\\
    &=30-90+60 &&\text{[\stepref{hcPointDefine}, \stepref{hcPointSecond},
                           \ref{hc:def:design}]}\\
    &=0 &&\text{[integer arithmetic].}
  \end{align*}\qed
 \end{proof}
 \step{hcPointZero}{For every $p\in P$, $(m(p)-2)^2=0$.}
 \begin{proof}\pf\ Each integer square is nonnegative; apply the
  zero-sum criterion of \ref{hc:facts:finite} to \stepref{hcPointSquares}.~\qed\end{proof}
 \qedstep
 \begin{proof}\pf\ An integer square vanishes only at zero, so
  \stepref{hcPointZero} gives $m(p)-2=0$ and $m(p)=2$ for every $p$.~\qed\end{proof}
\end{proof}

\subsubsection{The labelled extension and the matrix it supplies}

\begin{namedresult}{Definition}{(clique incidence and extended blocks)}\label{hc:def:extension}
 \lean{Design.sixCliqueIncMatrix}\lean{Design.hallConnorBlocks}
 \lean{Design.hallConnorIncMatrix}
 \uses{hc:def:design,hc:def:matrices,hc:def:graph}
 \leavevmode\par
 \assume{\begin{pfenum}
  \item $P,\cB$ are finite types with decidable equality.
  \item $D$ is a design with parameters $(15,5,21,7,2)$.
 \end{pfenum}}
 \pflet{$T\in\mat{\cB}{\cS}{\ints}$ have $T(B,S)=1$ for $B\in S$
 and $0$ otherwise. Put $P^*\triangleq P\sqcup\cS$ and
 $\cB^*\triangleq\operatorname{Option}(\cB)$.}
 Denote their injections by $\iota:P\to P^*$ and $\eta:\cS\to P^*$,
 the old labels by $\operatorname{some}(B)$, and the new label by $\operatorname{none}$.
 Define the block function $F:\cB^*\to\operatorname{Finset}(P^*)$ by
 \[
 \begin{split}
 F(\operatorname{some}(B))&\triangleq
    \iota(B)\cup\{\eta(S):S\in\cS,\ B\in S\},\\
 F(\operatorname{none})&\triangleq\{\eta(S):S\in\cS\}.
 \end{split}
 \]
 Here $\iota(B)$ is the image of the point set of $B$.
 Define $N^*\triangleq D.\mathrm{hallConnorIncMatrix}
 \in\mat{P^*}{\cB^*}{\ints}$ by the four entry formulas
 \[
 N^*(\iota(p),\operatorname{some}(B))\triangleq N(p,B),\quad N^*(\iota(p),\operatorname{none})\triangleq0,
 \]
 \[
 N^*(\eta(S),\operatorname{some}(B))\triangleq T(B,S),\quad N^*(\eta(S),\operatorname{none})\triangleq1.
 \]
\end{namedresult}

\begin{namedresult}{Lemma}{(parameters of the extended blocks)}\label{hc:lem:extension-parameters}
 \lean{Design.hallConnorBlocks_parameters}
 \uses{hc:def:design,hc:def:graph,hc:def:extension,hc:cor:two,hc:lem:seven,
 hc:lem:clique-meet,hc:lem:point-clique,hc:facts:finite}
 \leavevmode\par
 \assume{\begin{pfenum}
  \item $P,\cB$ are finite types with decidable equality.
  \item $D$ is a design with parameters $(15,5,21,7,2)$.
 \end{pfenum}}
 \prove{For the types and block function of Definition~\ref{hc:def:extension},
 \[
 \begin{gathered}
 |P^*|=22,\qquad|\cB^*|=22,\qquad
 \forall A\in\cB^*,\ |F(A)|=7,\\
 \forall x\in P^*,\ |\{A:x\in F(A)\}|=7,\\
 \forall x,y\in P^*,\ x\ne y\Rightarrow
                  |\{A:x,y\in F(A)\}|=2.
 \end{gathered}
 \]}
\end{namedresult}
\begin{proof}
 \pfsketch\ We will verify every cardinality and each of the three
 unordered types of point pair. The disjoint sum distinguishes new
 points from old points even if their underlying representations agree.
 \step{hcExtPoints}{$|P^*|=22$ and $|\cB^*|=22$.}
 \begin{proof}\pf\ Definition~\ref{hc:def:extension} gives the sums
  $|P|+|\cS|=15+7$ and $|\cB|+1=21+1$, using
  Definition~\ref{hc:def:design}, Lemma~\ref{hc:lem:seven}, and
  \path{Fintype.card_sum} and \path{Fintype.card_option} in
  \ref{hc:facts:finite}.~\qed\end{proof}
 \step{hcExtOldBlock}{For every $B\in\cB$, $|F(\operatorname{some}(B))|=7$.}
 \begin{proof}\pf\ There are five old points by
  Definition~\ref{hc:def:design} and two new points by
  Corollary~\ref{hc:cor:two}. Both $\iota$ and $\eta$ are injective,
  and their images are disjoint by Definition~\ref{hc:def:extension}.
  Apply \path{Finset.card_image_of_injective} and the disjoint-union
  count in \ref{hc:facts:finite} to get $5+2=7$.~\qed\end{proof}
 \step{hcExtNewBlock}{$|F(\operatorname{none})|=7$.}
 \begin{proof}\pf\ The map $\eta$ is injective, and
  Lemma~\ref{hc:lem:seven} gives seven members of $\cS$;
  use Definition~\ref{hc:def:extension} and
  \path{Finset.card_image_of_injective} in \ref{hc:facts:finite}.~\qed\end{proof}
 \step{hcExtOldReplication}{Each old point $\iota(p)$ lies in seven extended blocks.}
 \begin{proof}\pf\ Definition~\ref{hc:def:extension} preserves all
  incidences of $p$ with old labels and puts no old point in $F(\operatorname{none})$.
  Its old replication number is seven by Definition~\ref{hc:def:design}.~\qed\end{proof}
 \step{hcExtNewReplication}{Each new point $\eta(S)$ lies in seven extended blocks.}
 \begin{proof}\pf\ Definition~\ref{hc:def:extension} puts it in
  exactly the old labels $B\in S$, of which there are six by
  Definition~\ref{hc:def:graph}, and also in the distinct new
  label $\operatorname{none}$. The count is $6+1=7$.~\qed\end{proof}
 \step{hcExtOldOld}{\assume{$p,q\in P$ and $p\ne q$.}
  \prove{$\iota(p),\iota(q)$ lie together in exactly two extended blocks.}}
 \begin{proof}
  \step{hcExtOldOldOld}{The common old labels are precisely
   $\{\operatorname{some}(B):p,q\in B\}$.}
  \begin{proof}\pf\ The old-incidence clause of
   Definition~\ref{hc:def:extension}.~\qed\end{proof}
  \step{hcExtOldOldNew}{The new block contains neither old point.}
  \begin{proof}\pf\ The two injections have disjoint images in
   Definition~\ref{hc:def:extension}.~\qed\end{proof}
  \qedstep
  \begin{proof}\pf\ Step \stepref{hcExtOldOldOld} has cardinality two
   by Definition~\ref{hc:def:design} and \asm{1};
   \stepref{hcExtOldOldNew} excludes any further label.~\qed\end{proof}
 \end{proof}
 \step{hcExtOldNew}{\assume{$p\in P$ and $S\in\cS$.}
  \prove{$\iota(p),\eta(S)$ lie together in exactly two extended blocks.}}
 \begin{proof}
  \step{hcExtMixedOld}{The common old labels are precisely
   $\{\operatorname{some}(B):B\in S,\ p\in B\}$.}
  \begin{proof}\pf\ Use both old-label incidence clauses in
   Definition~\ref{hc:def:extension}.~\qed\end{proof}
  \step{hcExtMixedNew}{The new block does not contain $\iota(p)$.}
  \begin{proof}\pf\ Definition~\ref{hc:def:extension}.~\qed\end{proof}
  \qedstep
  \begin{proof}\pf\ Lemma~\ref{hc:lem:point-clique} gives two labels
   in \stepref{hcExtMixedOld}; \stepref{hcExtMixedNew} excludes
   a contribution from the new label.~\qed\end{proof}
 \end{proof}
 \step{hcExtNewOld}{\assume{$S\in\cS$ and $p\in P$.}
  \prove{$\eta(S),\iota(p)$ lie together in exactly two extended blocks.}}
 \begin{proof}\pf\ The set of common labels is unchanged when the
  two membership conditions are interchanged. Apply
  \stepref{hcExtOldNew} to $p,S$.~\qed\end{proof}
 \step{hcExtNewNew}{\assume{$S,T\in\cS$ and $S\ne T$.}
  \prove{$\eta(S),\eta(T)$ lie together in exactly two extended blocks.}}
 \begin{proof}
  \step{hcExtNewNewOld}{The common old labels are
   $\{\operatorname{some}(B):B\in S\cap T\}$, of size one.}
  \begin{proof}\pf\ Definition~\ref{hc:def:extension} gives this
   set, and Lemma~\ref{hc:lem:clique-meet} with \asm{1} gives its size.~\qed\end{proof}
  \step{hcExtNewNewNew}{Both points lie in the new block.}
  \begin{proof}\pf\ The clause for $F(\operatorname{none})$ in
   Definition~\ref{hc:def:extension}.~\qed\end{proof}
  \qedstep
  \begin{proof}\pf\ The old label of \stepref{hcExtNewNewOld}
   and the new label of \stepref{hcExtNewNewNew} are distinct,
   and together exhaust the common labels by
   Definition~\ref{hc:def:extension}.~\qed\end{proof}
 \end{proof}
 \qedstep
 \begin{proof}\pf\ Steps \stepref{hcExtPoints}--\stepref{hcExtNewReplication}
  verify the cardinalities, block sizes, and replication numbers.
  Steps \stepref{hcExtOldOld}, \stepref{hcExtOldNew},
  \stepref{hcExtNewOld}, \stepref{hcExtNewNew} exhaust ordered
  distinct point pairs: each point is in exactly one summand.
  Counting labels uses \path{Fintype.sum_option}, and splitting
  point sums uses \path{Fintype.sum_sum_type}, from
  \ref{hc:facts:finite}. These give every asserted equality.~\qed\end{proof}
\end{proof}

\begin{namedresult}{Definition}{(the Hall--Connor extension)}\label{hc:def:extension-design}
 \lean{Design.hallConnorExtension}
 \uses{hc:def:design,hc:def:extension,hc:lem:extension-parameters}
 \leavevmode\par
 \assume{\begin{pfenum}
  \item $P,\cB$ are finite types with decidable equality.
  \item $D$ is a design with parameters $(15,5,21,7,2)$.
 \end{pfenum}}
 Define $D^*\triangleq D.\mathrm{hallConnorExtension}$ of type
 \[
 \texttt{Design (Sum P D.sixCliques) (Option B) 22 7 22 7 2}
 \]
 with $D^*.\mathrm{block}\triangleq D.\mathrm{hallConnorBlocks}=F$
 from Definition~\ref{hc:def:extension}. The subtype of members
 is intended by \texttt{D.sixCliques} in the point type.
 The cardinality and incidence fields are exactly the equalities in
 Lemma~\ref{hc:lem:extension-parameters}; the positivity fields of
 Definition~\ref{hc:def:design} are $22>1$, $7>1$, and $22>0$.
\end{namedresult}

\begin{namedresult}{Theorem}{(the Hall--Connor residual identity)}\label{hc:thm:extension}
 \textup{[\cite{HC}, Theorems 2.1 and 3.2 at $k=5$.]}
 \lean{Design.hallConnorExtension_isSymmetric}
 \lean{Design.hallConnorExtension_residual}
 \uses{hc:def:design,hc:def:extension,hc:def:extension-design}
 \leavevmode\par
 \assume{\begin{pfenum}
  \item $P,\cB$ are finite types with decidable equality.
  \item $D$ is a design with parameters $(15,5,21,7,2)$.
 \end{pfenum}}
 \prove{The design $D^*=D.\mathrm{hallConnorExtension}$ is symmetric,
 and for every $B\in\cB$,
 \[
 D^*.\mathrm{block}(\operatorname{some}(B))\diff
 D^*.\mathrm{block}(\operatorname{none})=\iota(B).
 \]
 Thus deleting the new block and all its points recovers $D$ label by label.}
\end{namedresult}
\begin{proof}
 \pfsketch\ The parameter type gives symmetry; the two disjoint
 point images give the residual identity.
 \step{hcExtSymmetric}{$D^*$ is symmetric.}
 \begin{proof}\pf\ Its parameters in Definition~\ref{hc:def:extension-design}
  have $v=b=22$, which is Definition~\ref{hc:def:design}.~\qed\end{proof}
 \step{hcExtResidual}{For every $B$, the asserted block difference is $\iota(B)$.}
 \begin{proof}\pf\ Unfold the block function in
  Definition~\ref{hc:def:extension-design}. In
  Definition~\ref{hc:def:extension}, every added point belongs to
  $F(\operatorname{none})$ and no old point does. Removing that
  block from $F(\operatorname{some}(B))$ leaves exactly $\iota(B)$.~\qed\end{proof}
 \qedstep
 \begin{proof}\pf\ Steps \stepref{hcExtSymmetric} and \stepref{hcExtResidual}.~\qed\end{proof}
\end{proof}

\begin{namedresult}{Lemma}{(integral products for the extension)}\label{hc:lem:extension-products}
 \lean{Design.incMatrix_mul_sixCliqueIncMatrix}
 \lean{Design.sixCliqueIncMatrix_transpose_mul}
 \lean{Design.hallConnorIncMatrix_mul_transpose}
 \uses{hc:def:design,hc:def:matrices,hc:def:graph,hc:def:extension,hc:lem:incidence,
 hc:lem:clique-meet,hc:lem:point-clique,hc:facts:finite,hc:facts:matrix}
 \leavevmode\par
 \assume{\begin{pfenum}
  \item $P,\cB$ are finite types with decidable equality.
  \item $D$ is a design with parameters $(15,5,21,7,2)$.
 \end{pfenum}}
 \prove{All of the following are identities over $\ints$:
 \[
 NT=2J_{P,\cS},\qquad T^TT=5I_{\cS}+J_{\cS},\qquad
 N^*N^{*T}=5I_{P^*}+2J_{P^*}.
 \]}
\end{namedresult}
\begin{proof}
 \pfsketch\ We will expose the four matrix blocks so that the main
 contradiction can consume the integral identity directly.
 \step{hcProductsNT}{$NT=2J_{P,\cS}$.}
 \begin{proof}\pf\ By Definition~\ref{hc:def:extension} and
  \ref{hc:facts:matrix}, the $(p,S)$ entry counts $B\in S$ with
  $p\in B$. Lemma~\ref{hc:lem:point-clique} gives two.~\qed\end{proof}
 \step{hcProductsTTDiag}{$(T^TT)(S,S)=6$ for every $S\in\cS$.}
 \begin{proof}\pf\ This entry counts the six labels in $S$, by
  Definitions~\ref{hc:def:graph}, \ref{hc:def:extension}, and
  \ref{hc:facts:matrix}.~\qed\end{proof}
 \step{hcProductsTTOff}{\assume{$S,T_0\in\cS$ and $S\ne T_0$.}
  \prove{$(T^TT)(S,T_0)=1$.}}
 \begin{proof}\pf\ The entry counts $S\cap T_0$, by
  Definition~\ref{hc:def:extension} and \ref{hc:facts:matrix};
  Lemma~\ref{hc:lem:clique-meet} gives one.~\qed\end{proof}
 \step{hcProductsTT}{$T^TT=5I_{\cS}+J_{\cS}$.}
 \begin{proof}\pf\ The diagonal and off-diagonal cases
  \stepref{hcProductsTTDiag} and \stepref{hcProductsTTOff} are
  exhaustive and have the asserted entries $6$ and $1$.~\qed\end{proof}
 \step{hcProductsOO}{The old-old block of $N^*N^{*T}$ is $5I_P+2J_P$.}
 \begin{proof}\pf\ Sum over old labels and the new label using
  Definition~\ref{hc:def:extension}, \ref{hc:facts:matrix}, and
  \path{Fintype.sum_option} in \ref{hc:facts:finite}.
  The result is $NN^T+0$, equal to $5I_P+2J_P$ by
  Lemma~\ref{hc:lem:incidence}.~\qed\end{proof}
 \step{hcProductsON}{The old-new block of $N^*N^{*T}$ is $2J_{P,\cS}$.}
 \begin{proof}\pf\ The two label contributions are $NT$ and zero,
  by Definition~\ref{hc:def:extension}, \ref{hc:facts:matrix}, and
  \path{Fintype.sum_option} in \ref{hc:facts:finite};
  use \stepref{hcProductsNT}.~\qed\end{proof}
 \step{hcProductsNO}{The new-old block of $N^*N^{*T}$ is $2J_{\cS,P}$.}
 \begin{proof}\pf\ It is the transpose of the old-new block,
  since a matrix times its transpose is symmetric by
  \ref{hc:facts:matrix}; apply \stepref{hcProductsON}.~\qed\end{proof}
 \step{hcProductsNN}{The new-new block of $N^*N^{*T}$ is $5I_{\cS}+2J_{\cS}$.}
 \begin{proof}\pf\ Old labels contribute $T^TT$ and the new label
  contributes $J_{\cS}$, by Definition~\ref{hc:def:extension} and
  \ref{hc:facts:matrix}, using \path{Fintype.sum_option} in
  \ref{hc:facts:finite}. Step \stepref{hcProductsTT} gives their sum.~\qed\end{proof}
 \qedstep
 \begin{proof}\pf\ Steps \stepref{hcProductsOO}--\stepref{hcProductsNN}
  give the four blocks of $5I_{P^*}+2J_{P^*}$, because distinct
  summands never give a diagonal entry. Together with
  \stepref{hcProductsNT} and \stepref{hcProductsTT}, they prove all assertions.~\qed\end{proof}
\end{proof}

\begin{namedresult}{Theorem}{(Hall--Connor, matrix form)}\label{hc:thm:matrix}
 \lean{Design.hall_connor}
 \uses{hc:def:design,hc:def:extension,hc:lem:seven,hc:lem:extension-products,
 hc:facts:finite,hc:facts:matrix}
 \leavevmode\par
 \assume{\begin{pfenum}
  \item $P,\cB$ are finite types with decidable equality.
  \item $D$ is a design with parameters $(15,5,21,7,2)$.
 \end{pfenum}}
 \prove{There is $M\in\mat{\mathrm{Fin}\,22}{\mathrm{Fin}\,22}{\ints}$
 with $MM^T=5I_{22}+2J_{22}$.}
\end{namedresult}
\begin{proof}
 \pfsketch\ We will reindex both finite types in the integral
 identity, keeping the summed block index the same in both factors.
 \step{hcMatrixSizes}{$|P^*|=22$ and $|\cB^*|=22$.}
 \begin{proof}\pf\ Definition~\ref{hc:def:extension} gives disjoint
  sums of sizes $15+7$ and $21+1$, by
  Definition~\ref{hc:def:design}, Lemma~\ref{hc:lem:seven}, and
  \ref{hc:facts:finite}.~\qed\end{proof}
 \step{hcMatrixEquiv}{\pick\ equivalences
  $a:\mathrm{Fin}\,22\to P^*$ and $b:\mathrm{Fin}\,22\to\cB^*$.}
 \begin{proof}\pf\ Apply the enumeration fact in
  \ref{hc:facts:finite} to \stepref{hcMatrixSizes}.~\qed\end{proof}
 \step{hcMatrixDefine}{\pflet{$M(i,j)\triangleq N^*(a(i),b(j))$.}
  This is an integral matrix.}
 \begin{proof}\pf\ Every entry of $N^*$ is an integer by
  Definition~\ref{hc:def:extension}; composition preserves its scalar type.~\qed\end{proof}
 \step{hcMatrixProduct}{$(MM^T)(i,j)=(N^*N^{*T})(a(i),a(j))$ for all $i,j$.}
 \begin{proof}\pf\ Both factors use the bijection $b$ on the
  summed label index. Apply \path{Matrix.transpose_submatrix}
  and \path{Matrix.submatrix_mul_equiv} from \ref{hc:facts:matrix}
  to \stepref{hcMatrixDefine}, using \stepref{hcMatrixEquiv}.~\qed\end{proof}
 \step{hcMatrixTransport}{$MM^T=5I_{22}+2J_{22}$.}
 \begin{proof}\pf\ Lemma~\ref{hc:lem:extension-products} and
  \stepref{hcMatrixProduct} give entries
  $5\,1_{a(i)=a(j)}+2$. The injection $a$ in
  \stepref{hcMatrixEquiv} makes $a(i)=a(j)$ equivalent to $i=j$,
  so these are exactly the asserted entries.~\qed\end{proof}
 \qedstep
 \begin{proof}\pf\ The witness \stepref{hcMatrixDefine} and its
  equation \stepref{hcMatrixTransport} prove the existential conclusion.~\qed\end{proof}
\end{proof}

\subsubsection{The determinant obstruction and the conclusion}

\begin{namedresult}{Lemma}{(an odd power of five)}\label{hc:lem:odd-power}
 \lean{Computations.C19.not_sq_eq_49_mul_5_pow_odd}
 \lean{Computations.C19.not_sq_eq_49_mul_5_pow_21}
 \uses{hc:facts:prime}
 \leavevmode\par
 \assume{$m\in\nn$ and $d\in\ints$.}
 \prove{$d^2\ne49\cdot5^{2m+1}$. In particular $d^2\ne49\cdot5^{21}$.}
\end{namedresult}
\begin{proof}
 \pfsketch\ We will induct on $m$, universally quantifying the
 integer at each stage so that the induction hypothesis applies to $d/5$.
 \step{hcOddAll}{For every $m\in\nn$ and every $d\in\ints$,
  $d^2\ne49\cdot5^{2m+1}$.}
 \begin{proof}
  \step{hcOddBase}{\assume{$d\in\ints$ and $d^2=49\cdot5$.}
   \prove{\textsc{False}.}}
  \begin{proof}
   \step{hcOddBaseDivide}{\pick\ $e\in\ints$ with $d=5e$.}
   \begin{proof}\pf\ The equality in \asm{2} gives $5\mid d^2$;
    \ref{hc:facts:prime} gives $5\mid d$.~\qed\end{proof}
   \step{hcOddBaseCancel}{$5e^2=49$.}
   \begin{proof}\pf\ Substitute \stepref{hcOddBaseDivide} into
    \asm{2} and cancel the nonzero integer factor five.~\qed\end{proof}
   \qedstep
   \begin{proof}\pf\ Step \stepref{hcOddBaseCancel} would make
    $49$ divisible by five, whereas $49=5\cdot9+4$.~\qed\end{proof}
  \end{proof}
  \step{hcOddStep}{\assume{\begin{pfenum}
    \item $m\in\nn$ and $e^2\ne49\cdot5^{2m+1}$ for every $e\in\ints$.
    \item $d\in\ints$ and $d^2=49\cdot5^{2m+3}$.
   \end{pfenum}}
   \prove{\textsc{False}.}}
  \begin{proof}
   \step{hcOddStepDivide}{\pick\ $e\in\ints$ with $d=5e$.}
   \begin{proof}\pf\ The exponent $2m+3$ is positive, so
    \asm[2]{2} gives $5\mid d^2$. Apply
    \ref{hc:facts:prime}.~\qed\end{proof}
   \step{hcOddStepCancel}{$e^2=49\cdot5^{2m+1}$.}
   \begin{proof}\pf\ Substitute \stepref{hcOddStepDivide} into
    \asm[2]{2}, use $5^{2m+3}=25\cdot5^{2m+1}$, and cancel $25\ne0$.~\qed\end{proof}
   \qedstep
   \begin{proof}\pf\ Apply \asm[1]{2} to the integer $e$ in
    \stepref{hcOddStepCancel}.~\qed\end{proof}
  \end{proof}
  \qedstep
  \begin{proof}\pf\ Natural induction using \stepref{hcOddBase} and
   \stepref{hcOddStep}, since $2(m+1)+1=2m+3$.~\qed\end{proof}
 \end{proof}
 \qedstep
 \begin{proof}\pf\ Step \stepref{hcOddAll} proves the general assertion;
  its instance $m=10$ gives exponent $21$.~\qed\end{proof}
\end{proof}

\begin{namedresult}{Theorem}{(the Bruck--Ryser--Chowla instance)}\label{hc:thm:BRC}
 \lean{Matrix.bruck_ryser_chowla_twenty_two}
 \uses{hc:lem:det,hc:lem:odd-power,hc:facts:matrix,hc:facts:det}
 \leavevmode\par
 \assume{$M\in\mat{\mathrm{Fin}\,22}{\mathrm{Fin}\,22}{\ints}$.}
 \prove{$MM^T\ne5I_{22}+2J_{22}$.}
\end{namedresult}
\begin{proof}
 \pfsketch\ We will take determinants over $\rats$, retaining
 the integer determinant to apply the descent at five.
 \step{hcBRCBad}{\assume{$MM^T=5I_{22}+2J_{22}$.}\prove{\textsc{False}.}}
 \begin{proof}
  \step{hcBRCInteger}{Put $d\triangleq\det_{\ints}M$.
   Then $d\in\ints$ and $\det_{\rats}(M_{\rats})=(d:\rats)$.}
  \begin{proof}\pf\ The integer determinant is an integer by its
   scalar type; the cast identity is \path{Int.cast_det} in
   \ref{hc:facts:det}.~\qed\end{proof}
  \step{hcBRCSquare}{$(d:\rats)^2=49\cdot5^{21}$.}
  \begin{proof}
   \pf\ Cast \asm{1} to $\rats$ using \ref{hc:facts:matrix}.
   Since $5\ne0$ in $\rats$, the determinant lemma applies, giving
   \begin{align*}
    (d:\rats)^2&=\det(M_{\rats}M_{\rats}^T)
       &&\text{[\stepref{hcBRCInteger}, \ref{hc:facts:det}]}\\
     &=\det(5I_{22}+2J_{22}) &&\text{[\asm{1}]}\\
     &=5^{21}(5+22\cdot2) &&\text{[\ref{hc:lem:det}]}\\
     &=49\cdot5^{21} &&\text{[rational arithmetic].}
   \end{align*}\qed
  \end{proof}
  \step{hcBRCBack}{$d^2=49\cdot5^{21}$ in $\ints$.}
  \begin{proof}\pf\ The cast $\ints\to\rats$ is injective and
   preserves sums and products; apply it to \stepref{hcBRCSquare}.~\qed\end{proof}
  \qedstep
  \begin{proof}\pf\ Lemma~\ref{hc:lem:odd-power} contradicts
   \stepref{hcBRCBack}.~\qed\end{proof}
 \end{proof}
 \qedstep
 \begin{proof}\pf\ Step \stepref{hcBRCBad} proves the required negation.~\qed\end{proof}
\end{proof}

\begin{namedresult}{Theorem}{(the parameter contradiction)}\label{hc:thm:false}
 \lean{Design.false_of_parameters_15_5_2}
 \uses{hc:def:design,hc:thm:matrix,hc:thm:BRC}
 \leavevmode\par
 \assume{\begin{pfenum}
  \item $P,\cB$ are finite types with decidable equality.
  \item $D$ is a design with parameters $(15,5,21,7,2)$.
 \end{pfenum}}
 \prove{\textsc{False}.}
\end{namedresult}
\begin{proof}
 \pfsketch\ The matrix supplied by the embedding is forbidden by the obstruction.
 \step{hcFalseMatrix}{\pick\ an integer matrix $M$ of order $22$
  with $MM^T=5I_{22}+2J_{22}$.}
 \begin{proof}\pf\ Apply Theorem~\ref{hc:thm:matrix} to $D$.~\qed\end{proof}
 \qedstep
 \begin{proof}\pf\ Apply Theorem~\ref{hc:thm:BRC} to the witness
  and equation in \stepref{hcFalseMatrix}.~\qed\end{proof}
\end{proof}

\begin{namedresult}{Corollary}{(nonexistence of a $2$-$(15,5,2)$ design)}\label{hc:cor:main}
 \lean{Design.not_nonempty_15_5_2}\lean{Design.not_pair_counts_15_5_2}
 \uses{hc:def:design,hc:lem:parameters,hc:thm:false}
 \leavevmode\par
 \prove{The type
 \[\texttt{Design (Fin 15) (Fin 21) 15 5 21 7 2}\]
 is empty.
 Moreover, for every finite point type $P$ of size $15$ and every
 finite label type $\cB$, both with decidable equality, there is no
 function $F:\cB\to\operatorname{Finset}(P)$ whose blocks have size
 five and whose distinct point pairs lie in exactly two labels.}
\end{namedresult}
\begin{proof}
 \pfsketch\ We will state both the bundled and unbundled forms.
 \step{hcMainBundled}{The specified design type is empty.}
 \begin{proof}\pf\ Every inhabitant would give \textsc{False}
  by Theorem~\ref{hc:thm:false} with $P=\mathrm{Fin}\,15$,
  $\cB=\mathrm{Fin}\,21$.~\qed\end{proof}
 \step{hcMainUnbundled}{\assume{$P,\cB,F$ have the properties in the second assertion.}
  \prove{\textsc{False}.}}
 \begin{proof}
  \step{hcMainParameters}{\pick\ $D:\texttt{Design P B 15 5 21 7 2}$
   with $D.\mathrm{block}(B)=F(B)$ for every $B\in\cB$.}
  \begin{proof}\pf\ Apply Lemma~\ref{hc:lem:parameters} to \asm{1}.~\qed\end{proof}
  \qedstep
  \begin{proof}\pf\ Apply Theorem~\ref{hc:thm:false} to
   \stepref{hcMainParameters}.~\qed\end{proof}
 \end{proof}
 \qedstep
 \begin{proof}\pf\ Steps \stepref{hcMainBundled} and \stepref{hcMainUnbundled}.~\qed\end{proof}
\end{proof}

\begin{namedresult}{Definition}{(complementary labelled blocks)}\label{hc:def:compl}
 \lean{Design.compl}\uses{hc:def:design,hc:facts:finite}
 \leavevmode\par
 \assume{\begin{pfenum}
  \item $P,\cB$ are finite types with decidable equality.
  \item $D$ is a design with parameters $(v,k,b,r,\lambda)$.
  \item $2\le v-k$.
 \end{pfenum}}
 Define $D.\mathrm{compl}$ with the same point and label types by
 \[
 \begin{gathered}
 D.\mathrm{compl}:\texttt{Design P B v (v-k) b (b-r) (b+lam-2*r)},\\
 D.\mathrm{compl}.\mathrm{block}(B)\triangleq P\diff D(B).
 \end{gathered}
 \]
 The differences are natural differences. No injectivity of the
 block function is assumed. The proof below constructs its fields.
\end{namedresult}
\begin{proof}
 \pfsketch\ We will give the complement constructor its full
 parameter type. The bound on the new block size is exactly the
 condition needed by Definition~\ref{hc:def:design}.
 \step{hcComplGeneralSize}{Every complementary block has size $v-k$.}
 \begin{proof}\pf\ Its old block is a subset of $P$ with size $k$;
  subtract from $|P|=v$ using Definition~\ref{hc:def:design}
  and the finite-difference formula of \ref{hc:facts:finite}.~\qed\end{proof}
 \step{hcComplGeneralPoint}{Every point belongs to $b-r$ complementary labels.}
 \begin{proof}\pf\ The $r$ old labels through it form a subset of
  the $b$ labels. The other labels are exactly its complementary
  incidences, by Definition~\ref{hc:def:design} and \ref{hc:facts:finite}.~\qed\end{proof}
 \step{hcComplGeneralPair}{\new\ \assume{$p,q\in P$ and $p\ne q$.}
  \prove{Exactly $b+\lambda-2r$ complementary labels contain $p,q$.}}
 \begin{proof}
  \step{hcComplGeneralUnion}{Let $A_p,A_q$ be the sets of old
   incident labels of $p,q$. Then
   $|A_p\cup A_q|+\lambda=2r$.}
  \begin{proof}\pf\ Both sets have size $r$ and their intersection
   has size $\lambda$, by \asm{1} and Definition~\ref{hc:def:design}; use
   the union-intersection identity in \ref{hc:facts:finite}.~\qed\end{proof}
  \step{hcComplGeneralCount}{If $c$ counts labels containing
   neither point, then $c+2r=b+\lambda$.}
  \begin{proof}\pf\ The complement of $A_p\cup A_q$ among labels
   has size $c$, so $c+|A_p\cup A_q|=b$ by
   \ref{hc:facts:finite}. Add \stepref{hcComplGeneralUnion}.~\qed\end{proof}
  \qedstep
  \begin{proof}\pf\ Step \stepref{hcComplGeneralCount} gives
   $2r\le b+\lambda$ and $c=b+\lambda-2r$ in $\nn$.
   These are the common complementary labels by the block definition.
   This is the ordinary proof's inclusion-exclusion argument stated
   for general parameters, so the constructor has a precise type.~\qed\end{proof}
 \end{proof}
 \qedstep
 \begin{proof}\pf\ The types keep cardinalities $v,b$;
  steps \stepref{hcComplGeneralSize}, \stepref{hcComplGeneralPoint},
  and \stepref{hcComplGeneralPair} give the three incidence fields.
  The new block size is greater than one by \asm[3]{0}; $v>1,b>0$
  come from the old design in Definition~\ref{hc:def:design}.~\qed\end{proof}
\end{proof}

\begin{namedresult}{Corollary}{(the complementary parameters)}\label{hc:cor:complement}
 \lean{Design.false_of_parameters_15_10_9}\lean{Design.not_nonempty_15_10_9}
 \uses{hc:def:design,hc:def:compl,hc:thm:false}
 \leavevmode\par
 \assume{\begin{pfenum}
  \item $P,\cB$ are finite types with decidable equality.
  \item $D$ is a design with parameters $(15,10,21,14,9)$.
 \end{pfenum}}
 \prove{\textsc{False}. Thus
 \texttt{Design (Fin 15) (Fin 21) 15 10 21 14 9} is empty.}
\end{namedresult}
\begin{proof}
 \pfsketch\ Complementing every labelled block produces the forbidden parameters.
 \step{hcComplDesign}{Put $D'\triangleq D.\mathrm{compl}$.
  Then $D':\texttt{Design P B 15 5 21 7 2}$.}
 \begin{proof}\pf\ Apply the constructor of Definition~\ref{hc:def:compl}
  to \asm[2]{0}, with $2\le15-10=5$.
  Its parameters reduce to $15-10=5$, $21-14=7$, and
  $21+9-2\cdot14=2$. These are definitionally equal numerals in Lean.~\qed\end{proof}
 \qedstep
 \begin{proof}\pf\ Apply Theorem~\ref{hc:thm:false} to
  \stepref{hcComplDesign}. Since the argument applies to arbitrary
  finite $P,\cB$, it also excludes an inhabitant on
  $\mathrm{Fin}\,15$ and $\mathrm{Fin}\,21$.~\qed\end{proof}
\end{proof}

\subsubsection{Interfaces, provenance, and checks}

The statements above make the design-to-matrix bridge explicit:
\path{Design.hall_connor} returns the integer matrix equation that
\path{Matrix.bruck_ryser_chowla_twenty_two} negates. The constructor
\path{Design.hallConnorExtension} has the separately named theorem
\path{Design.hallConnorExtension_residual}, with point type \texttt{Sum P D.sixCliques} and label type
\texttt{Option B}. Here \path{D.sixCliques} denotes the subtype of
members when used as a type. The existential matrix theorem does
not need a general residual-design API.

No mathematical repair of the ordinary proof is needed.
We have made its moments, entry identities, principal-matrix bounds,
and clique incidence counts separately citable, put the determinant
lemma before its first use, and expanded the local and extension
cases. These are changes of proof structure. The ordinary
proof already proves the clique property before counting eight
edges, already defines the global family as all six-cliques, and
already uses rank to exclude $Q_1,Q_2$. Those repairs of the source
are retained here. Neither the unused rational determinant-sign
remark nor the unused source-local patterns are dependencies.
Three steps are tagged \new: using the rank-one determinant lemma,
taking the determinant of each certified inverse identity, and
stating the complement count for general parameters to give
\path{Design.compl} an exact interface. They retain the same
nonexistence argument.

All direct dependencies have local labels prefixed by
\path{hc:}; there are no deferred blueprint labels. The names in
\verb|\lean| are declarations of the Lean formalization, except for the
Mathlib API cited explicitly with paths. Mathlib
was inspected at commit
\par\noindent\path{082e2d37e8b0463410cdb532e111cd43d5a66174}\par\noindent
(4 September 2026), the version pinned in \path{lake-manifest.json}.
Its rank identities apply over $\rats$; the analytic theorem about
nonnegative determinants over an \path{RCLike} field is unnecessary.
The determinant obstruction uses only the even-order instance,
proved here without a quadratic-form local-global theorem.

The six displayed matrices are those of the ordinary proof. The exact
checks \path{blueprint/checks/hallconnor/check_proof.py}, with the $T(7)$
control \path{check_t7.py}, read them from this file and check all six
matrices, the two quadratic values, and every entry of both inverse
certificates; they are the checks of the ordinary proof in
\path{cgss/proofs/hallconnor/checks/}. The control
exercises all $210$ ordered adjacent pairs; no existing design can
exercise the hypothetical incidence structure. These checks do not
establish the quantified counting and clique arguments, which are
provided by the proofs above.
\end{mergedproofsection}
